%#!latex main.tex
\chapter{Introduction to the theory of $p$-adic period mappings.}
\label{chap:period-mappings}
\addtocontents{toc}{\protect\par\vskip2mm\hskip20mm(Geometry)\par\vskip5mm}
\minitoc
\newpage
\section{A survey of moduli and period mappings over $\C$.}\label{sec:1}

\begin{abst}
Review of moduli, Shimura varieties, 
Hodge structures, period mappings, period domains, Gauss-Manin connections.
\end{abst}
   
\subsection{Moduli spaces.}\label{sub:1.1}

\begin{para}\label{para:1.1.1} The term ``moduli'' was
 coined by Riemann to describe the continuous essential
 parameters of smooth complex compact curves of a given genus $g$. The use
 of this term has been extended to many classification problems in 
 analytic or algebraic
 geometry: typically, a classification problem consists of a discrete 
 part (the dimension, and other
 numerical invariants...) and a continuous part (the moduli). 

 One can distinguish between local and global moduli spaces: 

\begin{enumerate}
 \item local moduli are the parameters of the most general small (or
       infinitesimal) deformation of a given analytic manifold or
       algebraic variety. Let for instance $X$ be a compact analytic
       manifold.  Then a famous theorem of Kuranishi asserts that there
       is a local deformation $\pi: \udl{X}\rightarrow S$ with base
       point $s$ ($X={\udl{X}}_s$) such that any other deformation
       comes, locally around $s$, from $\pi$ by base change; moreover
       the tangent space of $S$ at $s$ is isomorphic to
       $\mathrm{H}^1(X,(\Omega^1_X)^\vee)$. For curves, the dimension of
       this space is $3g-3$ if $g\geq 2$ (the number of Riemann's
       complex moduli).
 \item Global moduli spaces, when they exist, provide the
       solution of the continuous part of moduli problems. For instance,
       Douady has constructed the global moduli space for closed analytic
       subspaces of a given compact complex manifold. Although global moduli
       spaces were already studied in the nineteenth century (in the case of
       curves), the concept was first put on firm foundations by Grothendieck
       in terms of representable functors.
\end{enumerate}

The existence of a global moduli space is often obstructed by the
presence of automorphisms of the objects under classification. One
possibility is to rule out these automorphisms by imposing some extra
structure (rigidification). When all automorphism groups are finite,
another possibility is to look for moduli orbifolds rather than moduli
spaces
\end{para}

\begin{para}\label{para:1.1.2}
 The most classical moduli problems concern complex elliptic
 curves. These are classified by their $j$-invariant. However the affine
 $j$-line ${\A}^1$ is not a fine moduli space: indeed, in any family of
 elliptic curves parametrized by the whole affine line, all the fibers
 are isomorphic.


 One can rigidify the problem by introducing a ``level $N$ structure''
 \index{levelNstructure@level $N$ structure} on 
 elliptic curves $E$, \ie fixing an identification $\Ker[N]_A\cong
 ({\Z}/N{\Z})^2$ such that the Weil pairing is given by
 $(\xi,\eta)\mapsto \exp(2\pi i(\xi\wedge\eta))$ (here $[N]_A$ denotes the
 multiplication by $N$ on the elliptic curve $A$).


For $N\geq 3$, this rules out automorphisms of $A$, and the
corresponding moduli functor is representable. For $N=3$, the moduli
space is ${\A}^1\setminus \{\mu^3=1\}$; the universal elliptic curve
with level $3$ structure is the Dixon elliptic curve
\begin{equation*}
 x^3+y^3+3\mu xy = 1 \text{ (in affine coordinates $x,y$)}
\end{equation*}
with $\Ker[3]_A=$ the set of flexes.


For $N=2$, the automorphism $[-1]_A$ remains; there is an elliptic curve
with level $2$ structure over ${\P}^1\setminus \{0,1,\infty\}$, namely
the Legendre elliptic curve
\begin{equation*}
 y^2=x(x-1)(x-\lambda) \text{ (in affine coordinates $x,y$)}
\end{equation*} 
with
$\Ker[2]_A=\{(0,0),(1,0),(\lambda, 0),(\infty,\infty)\}$, but it is not
universal: another one is given by $\lambda y^2=x(x-1)(x-\lambda)$.
\end{para}

\begin{para}\label{para:1.1.3}
 These moduli problems admit two natural generalizations: moduli
 problems for curves of higher genus, and moduli problems for polarized
 abelian varieties (polarizations ensure that the automorphism groups
 are finite).

 There is a moduli orbifold for smooth compact complex curves of genus
 $g>0$, denoted by $\mathcal{M}_g$, which can be described as
 follows. The set of complex structures on compact connected oriented
 surface $S_g$ of genus $g$ is in a natural way a contractible complex
 manifold of dimension $3g-3$ (\resp the Poincar\'e upper half plane
 $\mathfrak{h}$ if $g=1$), the Teichm\"uller space $\mathcal{T}_g$. The
 mapping class group $\Gamma_g$, \ie the component of connected
 components of the group of orientation-preserving self-homeomorphisms of
 $S_g$, acts on right on $\mathcal{T}_g$, and $\mathcal{M}_g$ is the
 quotient $\mathcal{T}_g/\Gamma_g$. For $g=1$, one recovers
 $\mathcal{M}_1=\mathfrak{h}/SL(2,{\Z})$.


There is a moduli orbifold for principally polarized abelian varieties of
dimension
$g>0$, denoted by $\mathcal{A}_g$, which can be described as follows. For any
abelian variety $A$, one has an exact sequence 
\begin{equation*}
 0\rightarrow L\rightarrow \Lie A\rightarrow A\rightarrow 0
\end{equation*}
and a principal polarization corresponds to a positive
definite hermitian form $H$ on $\Lie A$ whose imaginary part induces a
perfect alternate form on $L$. Choosing a $\C$-basis
$\omega_1,\ldots,\omega_g$ of $\Omega^1(A)\cong (\Lie A)^\vee$ and a
symplectic $\Z$-basis $\gamma_1^+,\gamma_1^-,\ldots,\gamma_g^+,
\gamma_g^-$ of $L$ gives rise to a period matrices 
$\Omega^+=(\int_{\gamma^+_j}\omega_i)$, $\Omega^-=(\int_{\gamma^-_j}\omega_i)$,
and $\tau=(\Omega^-)^{-1}\Omega^+$ does not depend on the
$\omega_i$'s.  By the properties of $H$, $\tau$ belongs to the Siegel
upper half space $\mathfrak{h}_g$ of symmetric matrices with positive
definite imaginary part (Riemann relations); it is well-defined up to
right multiplication by a matrix in $Sp(2g,{\Z})$ (which is the same as
the standard left action on $\mathfrak{h}_g$ of the transposed
matrix). This construction identifies $\mathcal{A}_g$ with
$\mathfrak{h}_g/Sp(2g,{\Z})$.


The construction which associates to a curve its jacobian variety with
the principal polarization given by the theta divisor gives rise to an
immersion $\mathcal{M}_g\hookrightarrow\mathcal{A}_g$, the Torelli map.
\end{para}

\begin{para}\label{para:1.1.4}
 Following G.~Shimura, one also studies refined moduli problems for
 ``decorated'' (principally polarized) abelian varieties, by prescribing
 in addition that the endomorphism algebra contains a given simple
 $\Q$-algebra, and imposing a level $N$ structure. This gives rise to
 moduli orbifolds (\resp moduli spaces if $N\geq 3$) of ``PEL type''.

 More precisely, let $B$ be a simple finite-dimensional $\Q$-algebra,
 with a positive involution $\ast$. Let $V$ be a $B$-module of finite
 type, endowed with an alternate $\Q$-bilinear form such that
 $\an{bv,w}=\an{v,b^\ast w}$.  One denotes by $G$ the algebraic $\Q$-group of
 $B$-linear symplectic similitudes of $V$.

 Let $\mathcal{B}$ be a maximal order in $B$ stable under $\ast$, and
let $L$ be a lattice in $V$, stable under $\mathcal{B}$ and autodual for
$\an{\;,\;}$ (this will represent the structure of $\mathrm{H}^1(A,{\Z})$ for
the abelian varieties $A$ under consideration). We choose in addition a
lagrangian subspace $F^1_0\subset V_{\C}$ stable under $B$ (which
will be $\Omega^1(A_0)\cong (\Lie A_0)^\vee$ for one particular abelian
variety $A_0$ under consideration).

There is a moduli orbifold (\resp moduli space if $N\geq 3$)
for principally $\ast$-polarized\footnote{a principal
$\ast$-polarization is a symmetric $\mathcal{B}$-linear isomorphism
between the abelian variety and its dual, endowed with the transposed
action of $\mathcal{B}$ twisted by $\ast$} abelian varieties $A$,
together with a given action $\iota: \mathcal{B}\rightarrow \End A$
such that $\det(\iota(b)| \Omega^1(A))=\det(b| F^1_0)$ for all
$b\in \mathcal{B}$ (Shimura type condition), and a given
$\mathcal{B}/N\mathcal{B}$-linear symplectic similitude
$A[N]=\Ker[N]_A\rightarrow \Hom(L,{\Z}/N {\Z})$.


It amounts to the same to consider abelian varieties $A$ up to isogeny,
with $B$-action (with Shimura type condition), together with a
$\Q$-homogeneous principal $\ast$-polarization and a class of $B$-linear
symplectic similitudes $\mathrm{H}^1(A,{\hat {\Z}}\otimes
\Q)\rightarrow {\hat {\Z}}\otimes V$ modulo the group $\{g \in
G({\hat {\Z}}\otimes {\Q}) \mid (g-1)({\hat {\Z}}\otimes L)\subset {\hat
{\Z}}\otimes L\}$.

These moduli orbifolds are called Shimura orbifolds of ``PEL
type''. \index{PELtype@PEL type} They are algebraic, and defined over
 the number field 
$E={\Q}[\tr(b| F^1_0)]$ (the so-called reflex field). Their
connected components are defined over finite abelian extensions of $E$,
\cf \cite{tds}.
\end{para}

\subsection{Period mappings.}\label{sub:1.2} 

\begin{para}\label{para:1.2.1}
 In the case of curves, the construction of period mappings belongs to
 the nineteenth
 century. The idea is to study the deformations of a curve by looking at
 the variation of its periods.  Let $f:{\udl{X}}\rightarrow S$ be an
 algebraic family of projective smooth (connected) curves of genus $g$
 over $\C$. Let $\tilde S$ be the universal covering of $S$. Then the
 homology groups $\mathrm{H}_1(-,{\Z})$ of the fibers form a constant
 local system on $\tilde S$. Let us choose a symplectic basis
 $\gamma_1^+,\gamma_1^-,\ldots ,\gamma_g^+, \gamma_g^-$. Choosing,
 locally, an auxiliary $\C$-basis $\omega_1,\ldots,\omega_g$ of
 $f_\ast\Omega^1_{\udl{X} /S}$ gives rise as above to a period matrices
 $\Omega^+=(\int_{\gamma^+_j}\omega_i)$, 
 $\Omega^-=(\int_{\gamma^-_j}\omega_i)$,  
 and $\mathcal{P}=(\Omega^-)^{-1}\Omega^+$ does not depend on the
 $\omega_i$'s and is well-defined on $\tilde S$.


 We get a commutative square 

 \begin{equation*}
  \begin{CD}
   \til{S} @>\mathcal{P}>> \mathfrak{h}_g\\
   @VVV @VVV\\
   S @>\mathcal{P}>> \mathfrak{h}_g/Sp(2g,\Z)
  \end{CD}
 \end{equation*}
 Applying this construction to the universal case
 $S=\mathcal{M}_g$ ($\mathcal{T}_g=\til{\mathcal{M}}_g$ in the orbifold
 sense), the bottom line of this commutative square is nothing but an
 analytic interpretation of the Torelli map.
\end{para}

\begin{para}\label{para:1.2.2}
 An important feature of the period mapping $\mathcal{P}$ is that it is
 a {\it quotient of solutions of a fuchsian differential system}, with
 rational exponents at infinity. This comes out as follows. The vector
 bundle $f_\ast\Omega^1_{\udl{X} /S}$ is a subbundle (of rank $g$) of a
 vector bundle $\mathcal{H}$ (of rank $2g$) whose fibers are the De Rham
 cohomology spaces $\mathrm{H}^1_{\mathrm{DR}}(\udl{X}_s)$. This bundle
 $\mathcal{H}$ carries an integrable connection, the Gauss-Manin
 connection
 \begin{equation*}
  \nabla_{\mathrm{GM}}:\;\mathcal{H}\rightarrow \Omega^1_S\otimes \mathcal{H}.
 \end{equation*}
 The local system of germs of horizontal analytic sections is
 $R^1f_\ast{\C}$ (with fibers
 $\mathrm{H}^1(\udl{X}_s,{\C})$). Completing, locally, the basis
 $\omega_1,\ldots,\omega_g$ of sections of $f_\ast\Omega^1_{\udl{X}/S}$
 into a basis $\omega_1,\ldots,\omega_g,\eta_1,\ldots,\eta_{g}$ of
 sections $\mathcal{H}$ allows to write the connection in form of a
 differential system, and a full {\it solution matrix} is given by the
 ``full $2g\times 2g$-period matrix''
 \index{full 2g 2g-period matrix@full $2g\times 2g$-period matrix} 
 $\begin{pmatrix} \Omega^+ &
  \Omega^- \cr N^+ & N^-  \end{pmatrix}$, while
 $\mathcal{P}=(\Omega^-)^{-1}\Omega^+$.


 In the case of the Legendre elliptic curve over ${\P}^1\setminus
 \{0,1,\infty\}$ (\cf I.\ref{sec-tale}), one finds that $\mathcal{P}$
 is given by the 
 quotient $\tau$ of two solutions of the hypergeometric differential
 equation with parameters $\bigl(\frac{1}{2},\frac{1}{2},1\bigr)$.
\end{para}

\begin{para}\label{para:1.2.3}
 The theory of period mappings has been generalized to any projective
 smooth morphism $f:\udl{X}\rightarrow S$ by P. Griffiths
 \cite{poioam3}.  We assume that $S$ and the fibers of $f$ are connected for
 simplicity.  The idea is the following: by abstracting the structure of
 the periods of the fibers of $f$, one arrives at the notion of
 polarized Hodge structure; one then construct a map from the universal
 covering $\tilde S$ to a classifying space $\mathcal{D}$ of polarized
 Hodge structures.

 For a fixed $n$, the real cohomology spaces
 $V_s^{\R}:= \mathrm{H}^n(\udl{X}_s,{\R})$) form a local system $V^{\R}$
 on $S$. We denote by $V^{\C}$ its complexification. The vector bundle
 $\mathcal{O}_S\otimes V^{\C}$ is the analytification of an algebraic
 vector bundle $\mathcal{H}$ on $S$, and $V^{\C}$ is the local system of
 germs of analytic solutions of an algebraic fuchsian integrable
 connection $\nabla_{\mathrm{GM}}$ on $\mathcal{H}$, the Gauss-Manin
 connection.

 On the other hand, the spaces $V_s^{\R}$ carry a natural algebraic
 representation of ${\C}^\times$ (Hodge structure), which amounts to a
 decomposition $V_s^{\C}=\bigoplus_{p+q=n,\;p\geq 0}
 \;V^{p,q},\;\;\overline{V^{p,q}}=V^{q,p}$. When $s$ varies, the
 filtrations
 \begin{equation*}
  F_s^p:=\bigoplus_{p'+q=n,\;p'\geq p} V^{p',q}
 \end{equation*}
 are the fibers of a decreasing filtration $F_s^p$ of $\mathcal{H}$ by
 vector subbundles (the Hodge filtration). The Gauss-Manin connection
 fails to preserve the Hodge filtration by just one notch:
 $\nabla_{\mathrm{GM}}(F^p)\subset \Omega^1_S\otimes F^{p-1}$.


 The last piece of data is the polarization: a bilinear form $Q$ on
 $V_s^{\R}$, $Q(v,v')=(-1)^nQ(v',v)$, such that $(-i)^nQ(v,\ovl{v})$ is
 hermitian on $V_s^{\C}$, definite (of sign $(-1)^p$) on each
 $V^{p,q}=F_s^p\cap \overline F_s^q$, and such that $F_s^p$ is the
 orthogonal of $F_s^{n-p-1}$ in $V_s^{\C}$.
 
 Let $G_{\R}$ denote the real group of similitudes of $Q$. The
 space of isotropic flags of type $0\subset \ldots \subset F_s^p\subset
 F_s^{p-1}\subset \ldots \subset V_s^{\C} \;$ is a homogeneous space
 \begin{equation*}
  \mathcal{D}^\vee = P\backslash G_{\C}
 \end{equation*}
 for a suitable parabolic subgroup $P$.  
 

 The flags satisfying the above positivity condition (on each $F_s^p\cap
 \ovl{F}_s^q$) are classified by an analytic open submanifold, the
 {\it period domain},
 \begin{equation*}
  \mathcal{D}\subset \mathcal{D}^\vee ;\;\; \mathcal{D}
   \cong K\backslash G_\R^{\mathrm{ad}}({\R})^{0}
 \end{equation*}
 \index{000GRad@$G_\R^{\mathrm{ad}}$}
 for a suitable compact subgroup $K\;$ ($\mathcal{D}^\vee$ is called the
 compact dual of $\mathcal{D}$).


Over $\tilde S$, the local system $V^{\C}$ becomes constant, and the
construction which attaches to any $s\in \tilde S$ the classifying point
of the corresponding flag $0\subset \ldots \subset F_s^p\subset
F_s^{p-1}\subset \ldots \subset V_s^{\C}\;$ gives rise to a holomorphic
mapping, the {\it period mapping}

 \begin{equation*}
  \tilde{S}\xrightarrow{\mathcal{P}} \mathcal{D}.
 \end{equation*}
 Via Pl\"ucker coordinates, $\mathcal{P}$ is again
 given by quotients of solutions of the Gauss-Manin connection
 $\nabla_{\mathrm{GM}}$.  The projective monodromy group of
 $\nabla_{\mathrm{GM}}$ is a 
 subgroup $\Gamma \subset G_{\R}^{\mathrm{ad}}$ (well-defined up to
 conjugation), and we get a commutative square

\begin{equation*}
 \begin{CD}
  \tilde{S}@>\mathcal{P}>> \mathcal{D} @. \subset \;\mathcal{D}^{\vee}\\
  @VVV @VVV @.\\
  S@>\mathcal{P}>> \mathcal{D}/\Gamma.
 \end{CD}
\end{equation*}
\end{para}

\begin{para}\label{para:1.2.4} In the simple case of a
 family of elliptic curves and $n=1$, $G_{\R}=GL_{2,\R}$, the flag space
 is just $\mathcal{D}^\vee={\P}^1_{\C}$, $K=PSO(2)$ (the isotropy group
 of $i\in {\P}^1_{\C}$ in $PSL_2({\R})=G_\R^{\mathrm{ad}}({\R})^0$), and
 $\mathcal{D}\cong K\backslash G_\R^{\mathrm{ad}}({\R})^0$ is the
 Poincar\'e upper half plane $\mathfrak{h}$.

 More generally, for a polarized abelian
 scheme, $G_{\R}$ is the group of symplectic similitudes, and
 $\mathcal{D}^\vee$ is the grassmannian of lagrangian subspaces of
 ${\C}^{2g}$, and $\mathcal{D}$ is the Siegel upper half space.

 \bigskip
 When one deals with refined moduli problems for abelian schemes with PEL
 decoration as in \ref{para:1.1.4}, one should work with the group $G_{\R}$ of
 $B$-{\it linear} symplectic similitudes of $V_{\R}$ (notation as in
 \ref{para:1.1.4}). In that case, the corresponding period domain
 $\mathcal{D}$ is a symmetric domain.

 More precisely, the datum of our $B$-stable lagrangian subspace
 $F^1_0\subset V_{\C}$ (attached to our particular $A_0$ with
 $\mathrm{H}^1(A_0,{\C})=V_{\C}$) amounts to the datum of an algebraic
 $\R$-homomorphism $h_0: {\C}^\times\rightarrow G_{\R}\subset
 GL(V_{\R})$ such that $h_0$ acts by the characters $z$ and $\ovl{z}$ on
 $V_{\C}$ ($F^1$ is the subspace where $h_0$ acts through $z$), and such
 that the symmetric form $\an{v,h(i) w}$ is positive definite on
 $V_{\R}$. For any abelian variety $A$ with PEL structure of the correct
 type, together with a fixed isomorphism $\mathrm{H}^1(A,{\C})\cong
 V_{\C}$, the corresponding subspace
 $F^1\mathrm{H}^1(A,{\C})=\Omega^1(A)\subset V_{\C}$ is an eigenspace
 for some conjugate of $h_0$. The set of conjugates of $h_0$ under
 $G({\R})$ is a finite union of copies of a symmetric domain
 $\mathcal{D}$.  \par For any component $S$ of the {\it Shimura
 orbifold}, we thus get a commutative square
 
 \begin{equation*}
  \begin{CD}
   \tilde{S}@>\mathcal{P}>> \mathcal{D} @. \subset\; \mathcal{D}^{\vee}\\
   @V\mathcal{Q}VV @VV\mathcal{Q}V @.\\
   S@>\mathcal{P}>> \mathcal{D}/\Gamma
  \end{CD}
 \end{equation*}
 \par\noindent
 where the horizontal maps are {\it isomorphisms}, $\mathcal{Q}$ denotes
 the quotient maps, and $\Gamma$ is a congruence subgroup of level $N$
 in the semi-simple group $G^{\mathrm{ad}}$ (a conjugate of the standard
 one).

 \bigskip
 The complex manifold $\tilde S$ admits the following
 modular description. Let $U$ be the oriented real Lie group
 $\Hom(L,(\R/\Z)^\times)$ (with $\mathcal{B}$-action and polarization),
 and let $\hat U$ be the associated infinitesimal Lie group. The
 automorphism group $J$ of the decorated $\hat U$ coincides with
 $G(\R)^0$. Then $\tilde S$ is a moduli space for {\it marked} decorated
 abelian varieties, where the marking is an isomorphism $\rho: \hat
 A\rightarrow \hat U$ of associated infinitesimal decorated oriented
 real Lie groups (note that, by taking duals of Lie algebras, $\rho$
 amounts to an $B$-linear symplectic isomorphism
 $\mathrm{H}^1(A,\R)\cong V^{\R}$). Note that the standard right action
 of $J$ on $\tilde S$ is obtained by functoriality.

\newpage
\section{Preliminaries on $p$-divisible groups.}\label{sec:2}

\begin{abst}
 Definitions and basic theorems about $p$-divisible groups,
 quasi-isogenies, liftings and deformations.
\end{abst}
\end{para}

\subsection{$p$-divisible groups and quasi-isogenies.}\label{sub:2.1}

\begin{para}\label{para:2.1.1} We begin with four definitions. Let $S$ be a scheme. Let
$\Lambda$ and $\Lambda'$ be commutative group schemes over $S$.  

A
homomorphism $f:\Lambda \rightarrow \Lambda'$ is called an {\it isogeny}
if it is an (f.p.p.f.) epimorphism with finite locally free kernel.


Let $p$ be a prime number. $\Lambda$ is a {\it $p$-divisible group}
(or Barsotti-Tate group) if $\displaystyle
\Lambda=\underset{n}{\liminj}\Ker [p^n]$ and for all $n$, the multiplication
$[p^n]$ by $p^n$ on $\Lambda$ is an isogeny.

\par\noindent If $S$ is connected, the rank of $\Ker [p^n]$ is then of
the form $p^{hn}$, where $h$ is an integer called the height of $\Lambda$.  

From the fact that $[p]$ is an isogeny, it follows that for 
$p$-divisible groups, $\Hom_S(\Lambda,\Lambda')$ is a torsion-free 
${\Z}_p$-module. 

A {\it quasi-isogeny} of $p$-divisible groups $\Lambda, \Lambda'$ is a
global section $\rho$ of the Zariski sheaf $ \udl{\Hom}_S(\Lambda,
\Lambda')\otimes_{\Z}{\Q}$ such that there exists locally an integer $n$
for which $p^n\rho$ is an isogeny. By abuse, one writes $\rho: \Lambda
\rightarrow \Lambda'$ as for homomorphisms. We denote by
$\qisog_S(\Lambda, \Lambda')$\index{000qisog@qisog} the ${\Q}_p$-space of
quasi-isogenies.


The objects of the category of {\it $p$-divisible groups on $S$ up
to isogeny} are the $p$-divisible groups over $S$; morphisms between
$\Lambda $ and $\Lambda'$ are global sections of $ \udl{\Hom}_S(\Lambda,
\Lambda')\otimes_{\Z}{\Q}$. This is a $\Q_p$-linear category.
\end{para}

\begin{para}\label{para:2.1.2} The most important example of a
 $p$-divisible group is $\Lambda =A[p^\infty]$, the (inductive system
 of) $p$-primary torsion of an abelian scheme $A$ over $S$
 (\cite{pg}); in this case, the height 
$h$ is twice the relative dimension of $A$. 
\end{para}

\begin{para}\label{para:2.1.3}
 When $p$ is locally nilpotent on $S$, any $p$-divisible group $\Lambda$
 over $S$ is formally smooth; the completion $\hat\Lambda$ of a
 $p$-divisible group $\Lambda$ along the zero section is a formal Lie
 group \cite[VI,3.1]{grothendieck74:_group_barsot_tate_dieud}.  However
 ${\hat\Lambda}$ is not necessarily 
 itself a $p$-divisible group; it is if and only if the separable rank of
 the fibers of $\Ker[p]$ is a locally constant function on $S$ (in which
 case $\Lambda$ is an extension
 \begin{equation*}
  1\rightarrow \hat{\Lambda}\rightarrow \Lambda\rightarrow
   \Lambda^{\mathrm{et}}\rightarrow 1 
 \end{equation*} 
 of an ind-etale $p$-divisible group $\Lambda^{\mathrm{et}}$ by the
 infinitesimal $p$-divisible group $\hat\Lambda$, \cf
 \cite[III,7.4]{grothendieck74:_group_barsot_tate_dieud}).
 An example of an ind-etale (\resp infinitesimal)
 $p$-divisible group of height $h=1\;$ is
 $\displaystyle {\Q}_p/{\Z}_p=\underset{n}{\liminj} {\Z}/p^n{\Z}$ (\resp
 $\hat {\G}_m$).
\end{para}

\subsection{Two theorems on $p$-divisible groups.}\label{sub:2.2} 

\begin{para}\label{para:2.2.1}
 A considerable amount of work has been
 devoted to $p$-divisible groups and their 
 applications to the $p$-adic study of abelian varieties and their local
 moduli. Among the pioneers, let us mention: Barsotti, Tate, Serre,
 Grothendieck, Lubin, Messing...


 In the sequel, we shall outline the main concepts and results of the
 Rapoport-Zink theory of period mappings for $p$-divisible groups,
 emphasizing the analogy with the complex case. We refer to their book
 \cite{psfpg} for details and proofs. We shall also emphasize the
 relation to differential equations, which does not appear in
 \cite{psfpg}.

 The whole theory relies upon three basic theorems on $p$-divisible
 groups and their Dieudonn\'e modules, which we shall recall below:


 \begin{enumerate}
  \item the rigidity theorem for $p$-divisible groups up to isogeny,
  \item the Serre-Tate theorem, and
  \item the Grothendieck-Messing theorem,
 \end{enumerate}

 \noindent
 which deal with infinitesimal deformations of $p$-divisible groups.
\end{para}

\begin{para}\label{para:2.2.2} Let $\mathfrak{v}$ be a complete discrete
valued ring $\mathfrak{v}$ of mixed characteristic $(0,p)$. Let 
$\Nil_\mathfrak{v}$ denote the category of locally noetherian 
$\mathfrak{v}$-schemes $S$ on which $p$ is locally nilpotent.  For any
$S$ in $\Nil_\mathfrak{v}$, the ideal of definition 
$\mathcal{J}=\mathcal{J}_{S_{\mathrm{red}}}$ of the closed subscheme $S_{\mathrm{red}}$ 
(of characteristic $p$) is locally nilpotent.
\end{para}

\begin{thm}[Rigidity theorem for $p$-divisible groups up to isogeny]
 \label{thm:2.2.3}
 
 Every homomorphism
 $\;\ovl{\rho}:\;\Lambda\times_S S_{\mathrm{red}} \rightarrow \Lambda'\times_S
 S_{\mathrm{red}}$ of $p$-divisible groups up to isogeny admits a unique lifting
 $\rho: \Lambda \rightarrow \Lambda'$. Moreover, $\rho$ is a
 quasi-isogeny if $\ovl{\rho}$ is.
\end{thm}

\begin{proof}
 (following an argument of V. Drinfeld). We may assume that $S$ is
 affine, and that $\mathcal{J}$ is nilpotent. Then the connected part
 ${\hat\Lambda'}$ is a formal Lie group, hence there is a power $p^n$ of
 $p$ such that for any affine $S$-scheme $S'$, ${\hat\Lambda}'(S')$ is
 killed by $[p^n]$ (in fact if $\mathcal{J}^{r+1}=0$ on $S$, one can take
 $n=r^2$).

 We observe that $\Ker(\Lambda'(S')\rightarrow
 \Lambda'(S'_{\mathrm{red}}))= {\hat\Lambda}'(S')$. We have
 \begin{align*}
 &\Ker(\Hom(\Lambda(S') ,
 \Lambda'(S'))\rightarrow \Hom(\Lambda(
 S'_{\mathrm{red}}) , \Lambda'(S'_{\mathrm{red}}) )\\
 =&
 \Hom (\Lambda(S') ,\Ker(\Lambda'(S')
 \rightarrow \Lambda'(S'_{\mathrm{red}}))) 
 \end{align*}
 which is zero because $\Ker(\Lambda'(S')\rightarrow
 \Lambda'(S'_{\mathrm{red}})$ is killed by $[p^n]$ while $\Lambda$ is
 $p$-divisible.

 \noindent This implies the injectivity of
 $\qisog_S(\Lambda , \Lambda') \rightarrow 
 \qisog_{S_{\mathrm{red}}}(\Lambda\times_S S_{\mathrm{red}} ,
 \Lambda'\times_S S_{\mathrm{red}}) $ (taking into account the fact that 
 $\Hom(\Lambda( S'_{\mathrm{red}}), \Lambda'(S'_{\mathrm{red}}) )$ is
 torsion-free).

 \noindent For the surjectivity, let us first show that for any
 ${\ovl{f}}\in \Hom(\Lambda(S'_{\mathrm{red}}),
 \Lambda'(S'_{\mathrm{red}}))$, there is 
 a lifting $g\in \Hom(\Lambda(S'), \Lambda'(S'))$ of $p^n {\ovl{f}}$. For
 any $x\in \Lambda( S')$, let $\ovl{x}$ denote its image in $\Lambda(
 S'_{\mathrm{red}})$. Since $\Lambda'$ is formally smooth,
 ${\ovl{f}}({\ovl{x}})$ 
 admits a lifting $y\in \Lambda'(S')$.  Since
 $\Ker(\Lambda'(S')\rightarrow \Lambda'(S'_{\mathrm{red}}))) $ is killed
 by $p^n$, 
 $g(x):=[p^n] y$ is well-defined. This construction provides the desired
 lifting $g$.

 It remains to show that $g$ is an isogeny if
 $\ovl{f}$ is. Let $\ovl{f}'$ be a quasi-inverse of $\ovl{f}$
 ($\ovl{f}'\ovl{f}=[p^m]$), and $g'$ be a lifting of $p^n\ovl{f}'$. Then
 $g'g=[p^{2n+m}]$ by unicity of liftings. Thus $g$ is an epimorphism, and
 the subscheme $\Ker g$ of $\Ker [p^{2n+m}]$ is finite over $S$. On the
 other hand, $\Lambda$ is flat over $S$ and the fibers of $g'$ are flat
 (being isogenies). The criterium of flatness fiber by fiber
 \cite[11.3.10]{ega44} implies that $g'$ is flat, and we conclude that
 it is an isogeny.
\end{proof}

Let us observe that there is variant of both the statement and the
proof of \ref{thm:2.2.3} for abelian schemes, instead of $p$-divisible groups. 

 The next theorem asserts that deforming an abelian scheme of
 characteristic 
 $p$ is equivalent to deforming its $p$-divisible group. For any $S $ in
 $\Nil_\mathfrak{v}$, let $\Def_{p\text{-div}}(S)$ be the category of
 triples $(\ovl{A},\Lambda, \epsilon)$ consisting of an abelian scheme
 over $S_{\mathrm{red}}$, a $p$-divisible group $\Lambda$ over $S$ and an
 isomorphism $\epsilon: \ovl{A}[p^\infty]\cong \Lambda\times_S S_{\mathrm{red}}$.

\begin{thm}[The Serre-Tate theorem]\label{thm:2.2.4}
 The functor
 \begin{equation*}
  A\;\mapsto \;(A_{\mathrm{red}},\;\Lambda=A[p^\infty],
   \;\epsilon: A_{\mathrm{red}}[p^\infty]\cong \Lambda\times_S S_{\mathrm{red}})
 \end{equation*}
 induces an equivalence of categories between the
 category of abelian schemes over $S$ and the category
 $\Def_{p\text{-div}}(S)$.
\end{thm}


We again follow Drinfeld's argument \cite{copsr}, \cf also \cite{slm}.
We may assume that $S$ is affine and that
$\mathcal{J}^{n+1}=0$. We begin with the full faithfulness: given a
homomorphism $f^\infty: A[p^\infty]\rightarrow B[p^\infty]$ such that
$f^\infty$ $\pmod{\mathcal{J}}$ comes from a homomorphism $\ovl{f} :
A_{\mathrm{red}}\rightarrow B_{\mathrm{red}}$ of abelian schemes over $S_{\mathrm{red}}$, we have
to show that there exists a unique homomorphism $f: A\rightarrow B$
compatible with $f^\infty$ and $\ovl{f}$. The unicity follows from the
injectivity of $\Hom(A,B)\rightarrow \Hom(A_{\mathrm{red}},B_{\mathrm{red}})\;$
(\ref{thm:2.2.3}). Furthermore, \ref{thm:2.2.3} shows the existence of a
homomorphism $g$ 
compatible with $\;p^nf^\infty$ and $p^n\ovl{f}$. To show the existence
of an $f$ such that $g=p^nf=f\circ [p^n]$, the point is to show that $g$
kills $A[p^n]$, which can be seen on $g^\infty$.


For the essential surjectivity, we have to construct an abelian scheme
$A$ over $S$ compatible with a given datum $(\ovl{A},\Lambda, \epsilon)$
in $\Def_{p\text{-div}}(S)$.  By unicity of liftings, it suffices to do so
locally on $S$.  hence we may assume that $S$ is affine, and choose an
abelian scheme $A'$ over $S$ and an isogeny $A'_{\mathrm{red}}\rightarrow
\ovl{A}$. By the above argument, $p^n$ times the given isogeny
$A'_{\mathrm{red}}\rightarrow \ovl{A}$ lifts to a (unique) isogeny
$e:\;A'[p^\infty ]\rightarrow \Lambda$ of $p$-divisible groups over
$S$. Its kernel is a finite locally free subgroup scheme of $A'$, and we
can form the abelian scheme $A'':=A'/\Ker e$. The isogeny $A''[p^\infty
]\rightarrow \Lambda$ induced by $e$ is an isomorphism
$\pmod{\mathcal{J}}$, hence is an isomorphism.

\begin{para}\label{para:2.2.5} In the sequel, we shall have to deal
 occasionally with $p$-divisible groups over formal schemes rather 
than over schemes. Our formal schemes $\mathfrak{X}$ will be adic,
locally noetherian; hence there is a largest ideal of definition
$\mathcal{J}\subset \mathcal{O}_{\mathfrak{X}}$, and
$\mathfrak{X}=\liminj \;{\mathfrak{X}}_n\;,\;{\mathfrak{X}}_n =
\Spec(\mathcal{O}_{\mathfrak{X}}/\mathcal{J}^{n+1})\;\;
({\mathfrak{X}}_{\mathrm{red}}={\mathfrak{X}}_0)\;$.

A $p$-divisible group $\Lambda$ over $\mathfrak{X}$ will be an inductive
system of $p$-divisible groups $\Lambda_n$ over $\mathfrak{X}_n$ such
that $\Lambda_{n+1}\times_{{\mathfrak{X}}_{n+1}}{\mathfrak{X}}_n\cong
\Lambda_n$.

\newpage
\section{A stroll in the crystalline world.}\label{sec:3}

\begin{abst}
 Review of Dieudonn\'e modules, crystals, Grothendieck-Messing theory,
 convergent isocrystals.
\end{abst}
\end{para}

\subsection{From Dieudonn\'e modules to crystals.}\label{sub:3.1}

\begin{para}\label{para:3.1.1}
 At first, Dieudonn\'e theory presents itself as an analogue of Lie
 theory for formal Lie groups over a perfect field $k$ of characteristic
 $p>0$. In practice, its scope is limited to the case of {\it
 commutative} formal Lie groups (in which case the classical counterpart
 is not very substantial). The theory extends to the case of
 $p$-divisible groups over $k$.


 Let $W$ be the Witt ring of $k$, with its Frobenius automorphism
 $\sigma$. The theory associates to any 
 $p$-divisible group $\Lambda$ over $k$ a {\it Dieudonn\'e module}, \ie a
 finitely generated free module over $W$ 
 endowed with a $\sigma$-linear Frobenius $F$ and a $\sigma^{-1}$-linear
 Verschiebung 
 $V$ satisfying $FV=VF=p.\id$ (we work with the contravariant theory: 
 ${\D}(\Lambda)=\Hom(\Lambda , \mathcal{W}_k)$, 
 \index{000D@$\D$}
 where $\mathcal{W}_k\;$
 is the Witt scheme of $k$). This provides an 
 {\it anti-equivalence of categories between $p$-divisible groups over
 $k$ and Dieudonn\'e modules}. The rank of 
 ${\D}(\Lambda)$ is the height $h$ of $\Lambda$. 
\end{para}

\begin{exs}
  $\bullet \;\;{\D}(\Q_p/\Z_p)=W,$ with $F=\sigma,
 \;V=p\sigma^{-1}$. 
 \par\noindent $\bullet\;\; {\D}(\hat\G_m)=W(-1)$, \ie $W$ as an
 underlying module, with $F=p\sigma, 
 \;V=\sigma^{-1}$.
 \par\noindent $\bullet$ When
 $\Lambda=A[p^\infty]$ for an abelian variety $A$ over $k$,
 ${\D}(\Lambda)$ may be identified with 
 $\mathrm{H}^1_{\mathrm{cris}}(A/W)$. 
\end{exs}

 However, the structure of Dieudonn\'e modules is rather complicated in
 general unless one inverts $p$. 

\begin{para}\label{para:3.1.2}
 When $k$ is algebraically closed, Dieudonn\'e modules $\otimes \Q$ were
 completely classified by Dieudonn\'e (this provides a classification of
 $p$-divisible groups over $k$ up to isogeny). They are classified by
 their {\it slopes} $\lambda\in [0,1]\cap \Q$ and multiplicities
 $m_\lambda\in \Z_{>0}$ (these data are better recorded in the form of a
 Newton polygon). This comes as follows: each Dieudonn\'e module $\D$
 has a canonical increasing finite filtration by Dieudonn\'e submodules
 $\D_\lambda$, such that the associated graded $W$-module is free, and
 with the following property: if $\lambda=a/b$ in irreducible form,
 $\Gr_\lambda\otimes \Q$ admits a basis of $m_\lambda$ vectors $x$
 satisfying $F^bx=p^ax$.


 For instance, the slope of $W(-1)$ is $1$; the slope of the Dieudonn\'e
 module attached to 
 (the $p$-divisible group of) a supersingular elliptic curve over $k$ is
 $1/2$, with multiplicity $2$.
\end{para}

\begin{para}\label{para:3.1.3} The problem of generalizing Dieudonn\'e
 theory to $p$-divisible groups over more general bases $S$ (over which
$p$ is nilpotent) has been tackled and advertised by Grothendieck
 \cite{grothendieck74:_group_barsot_tate_dieud}, and further studied by many geometers.


 What should be the right substitutes for the $W$-modules ${\D}(\Lambda)$
 (\resp for the $W[\frac{1}{p}]$-spaces ${\D}(\Lambda)_{\Q}$)?


 Grothendieck's proposal was to define ${\D}(\Lambda)$ as a {\it
 $F$-crystal} on the crystalline site of $S$. 

 The right substitute for $W[\frac{1}{p}]$-spaces was proposed later by
 P. Berthelot and A. Ogus under the name of {\it convergent
 $F$-isocrystal}. Let us describe these notions.
\end{para}


\subsection{The Dieudonn\'e crystal of a $p$-divisible group.}\label{sub:3.2}

\begin{para}\label{para:3.2.1} Let $\mathfrak{v}$ be a finite extension of
 $W$ of degree $e$. We assume that $e<p$, and that $\varpi$ is a 
 uniformizing parameter of $\mathfrak{v}$ such that $\sigma$ extends to 
 $\mathfrak{v}$ by setting $\sigma(\varpi)=\varpi$.


Let $S$ be in $\Nil_\mathfrak{v}$. An
$S$-divided power thickening
$T_0\hookrightarrow T$ is given by
\begin{enumerate}
 \item an $S$-scheme $T_0$,
 \item a $\mathfrak{v}$-scheme $T$ on which $p$ is locally nilpotent,
 \item a closed immersion $T_0\hookrightarrow T$ over $\mathfrak{v}$,
 \item a collection of maps ``$\frac{x^n}{n!}$'' from the ideal of
 definition of $T_0$ in $T$ to $\mathcal{O}_T$, which 
 satisfy the formal properties of the divided powers. 
\end{enumerate}

 $S$-divided power thickenings form a category
 $(S/\mathfrak{v})_{\mathrm{cris}}$ (a fppf site). A sheaf on
 $(S/\mathfrak{v})_{\mathrm{cris}}$ is the data, for every
 $T_0\hookrightarrow T$, of a fppf sheaf $\mathcal{E}_{T_0,T}$ on $T$,
 and for every morphism $f: (T_0\hookrightarrow T)\rightarrow
 (T_0'\hookrightarrow T')$ in $(S/\mathfrak{v})_{\mathrm{cris}}$, of a
 homomorphism $f^\ast\mathcal{E}_{T_0,T}\rightarrow
 \mathcal{E}_{T_0',T'}$ satisfying the obvious transitivity
 condition. Example: $\mathcal{O}_{S_0/\mathfrak{v}}$, defined by
 $(\mathcal{O}_{S/\mathfrak{v}})_{T_0,T}=\mathcal{O}_T$. Another example:
 any fppf sheaf $\mathcal{F}$ on $S_0$ gives rise to a sheaf
 $\mathcal{F}_{\mathrm{cris}}$ on $(S/\mathfrak{v})_{\mathrm{cris}}$ by
 $(\mathcal{F}_{\mathrm{cris}})_{T_0,T}=\mathcal{F}_{T_0}$.

 A {\it crystal} on $(S/\mathfrak{v})_{\mathrm{cris}}$ (or,
 abusively, on $S$) is a sheaf $\mathcal{E}$ of
 $\mathcal{O}_{S/\mathfrak{v}}$-modules on
 $(S/\mathfrak{v})_{\mathrm{cris}}$ such that the homomorphisms
 $f^\ast\mathcal{F}_{T_0,T}\rightarrow \mathcal{F}_{T_0',T'}$ are
 isomorphisms (in Grothendieck's terms: ``crystals grow and are
 rigid''). One denotes by $\mathcal{E}_T$ the $\mathcal{O}_T$-module
 obtained by evaluating the crystal on $(\id: T\hookrightarrow T) \in
 (S/\mathfrak{v})_{\mathrm{cris}}$.

 Because of this rigidity, and using the canonical divided powers on the
 ideal $\varpi\mathcal{O}_S$ (since $e<p,\;\frac{\varpi^n}{n!}\in
 \mathfrak{v}$ for any $n$), one gets an equivalence of categories
 between crystals on $S$ and crystals on the scheme $S_0=S/\varpi$ of
 characteristic $p$.
\end{para}

\begin{para}\label{para:3.2.2}
 The Dieudonn\'e crystal attached to a $p$-divisible group $\Lambda$
 over $S$ is the crystalline $\mathcal{E}xt$-sheaf
 $\mathcal{E}xt^1(\Lambda_{\mathrm{cris}},
 \mathcal{O}_{S/\mathfrak{v}})$. It can be shown that ${\D}(\Lambda)$ is
 a finite locally free crystal on $S$, of rank the height of
 $\Lambda$. It depends functorially in $\Lambda$ (in a contravariant
 way).

 An important, though formal, consequence of the crystalline local
 character of this definition is that {\it the formation of
 ${\D}(\Lambda)$ commutes with base change in}
 $(S/\mathfrak{v})_{\mathrm{cris}}$, and for any $S$-divided power
 thickening $(T_0\hookrightarrow T)$, there is a canonical isomorphism
 of $\mathcal{O}_T$-modules ${\D}(\Lambda)_{T_0,T}\cong
 {\D}(\Lambda_{T_0})_{T_0,T}.$

 It follows that the datum of ${\D}(\Lambda)$ amounts to the datum of
 ${\D}(\Lambda\times_S S_0)$: the Dieudonn\'e crystal depends only on the
 $p$-divisible group modulo $\varpi$.
\end{para}

\begin{para}\label{para:3.2.3}
 The pull-back of a crystal $\mathcal{E}$ on $S_0$ via the absolute
 Frobenius $F_{S_0}$ is denoted by $\mathcal{E}^{(p)}$. An $F$-crystal
 structure on $\mathcal{E}$ is the datum of a morphism of crystals
 $F:\;\mathcal{E}^{(p)}\rightarrow \mathcal{E}$ which admits an inverse
 up to a power of $p$.

 If $\Lambda$ is a $p$-divisible group over $S_0$, ${\D}(\Lambda)$ is
 endowed by functoriality with a structure of $F$-crystal. The classical
 Dieudonn\'e module is obtained for
 $S_0=T_0=\Spec(k),\;T=\Spf(W)=\liminj\;\Spec(W/p^n)$.


 Although we shall not use this fact, let us mention that the
 ``Dieudonn\'e functor'' ${\D}$ from $p$-divisible groups on $S_0$ to
 $F$-crystals is fully faithful under further mild assumptions on $S_0$
 (and even an equivalence with values in the category of ``Dieudonn\'e
 crystals'', \ie finite locally free $F$-crystals for which $F$ admits an
 inverse up to $p$ --- the Verschiebung); without any further assumption on
 $S_0$, it induces an equivalence of categories between ind-etale
 $p$-divisible groups on $S_0$ and unit-root $F$-crystals (\ie finite
 locally free $F$-crystals for which $F$ is an isomorphism), 
 \cf \cite{tddc1}, \cite{jong95:_cryst_dieud}.
\end{para}


\subsection{The Hodge filtration.}\label{sub:3.3}

\begin{para}\label{para:3.3.1}
 Let $\Lambda^\vee$ be the (Serre) dual of $\Lambda$ (the $\liminj$ of
 the Cartier duals of the $\Lambda[p^n]$'s). There is a canonical,
 locally split, exact sequence of $\mathcal{O}_S$-modules
 \begin{equation*}
  0\rightarrow F^1\rightarrow {\D}(\Lambda)_S\rightarrow
   \Lie \Lambda^\vee\rightarrow 0
 \end{equation*}
 where $F^1=\CoLie \Lambda=\omega_\Lambda$ may be identified with the
 module of invariant differentials on the formal Lie group
 $\hat\Lambda$. In case $\Lambda$ is the $p$-divisible group attached to
 an abelian $S$-scheme $\udl{A}$, this exact sequence reduces to the
 Hodge exact sequence
 \begin{equation*}
  0\rightarrow F^1\rightarrow \mathrm{H}^1_{\mathrm{DR}}(\udl{A}/S)\rightarrow
   \Lie \udl{A}^\vee\rightarrow 0.
 \end{equation*}
\end{para}

\begin{para}\label{para:3.3.2}
 Now let $S_0\hookrightarrow S$ be a {\it nilpotent} divided power
 thickening (this means that the products
 ``$\frac{x_1^{n_1}}{{n_1}!}$''$\ldots$``$\frac{x_1^{n_m}}{{n_m}!}$'' vanish
 for $n_1+\ldots+n_m\gg 0$); this is for instance the case if
 $S_0=S/\varpi$ as before and $e<p-1$.


 The next theorem asserts that deformations of $p$-divisible
 groups are controlled by the variation of the Hodge filtration in the
 Dieudonn\'e module \cite{messing72:_barsot_tate}. For any $p$-divisible group
 $\ovl{\Lambda}$ on $S_0$, and any lifting $\Lambda$ over $S$, we have
 seen that ${\D}(\Lambda)_S$ identifies with ${\D}(\ovl{\Lambda})_S$,
 which is a lifting of ${\D}(\ovl{\Lambda})_{S_0}$ to $S$.
\end{para}

\begin{thm}[The Grothendieck-Messing theorem]\label{thm:3.3.3}
 The functor
 \begin{equation*}
  \Lambda\;\mapsto (\Lambda\times_S S_0,\;\omega_\Lambda)
 \end{equation*}
 induces an equivalence of categories between the category of
 $p$-divisible groups on $S$, and the category of pairs consisting of a
 $p$-divisible group $\ovl{\Lambda}$ over $S_0$ (Zariski-locally liftable
 to $S$) and a locally direct factor vector subbundle of
 ${\D}(\ovl{\Lambda})_S$ which lifts $\;\omega_{\ovl{\Lambda}}\subset
{\D}(\ovl{\Lambda})_{S_0}$.
\end{thm}


 We have already seen the faithfulness in \ref{thm:2.2.3}.  We just
 indicate here 
 the principle of proof of essential surjectivity and fullness. Working
 with nilpotent divided power thickenings allows one to construct, by
 the exponential method, a formally smooth group scheme
 $\;\E(\ovl{\Lambda})\;$ over $\;S\;$ together with a canonical
 isomorphism
 ${\D}(\ovl{\Lambda})_{S_0}=\Lie (\E(\ovl{\Lambda})\times_S{S_0})$, in
 such a way that the canonical exact sequence

 \begin{equation*}
  0\rightarrow \omega_{\ovl{\Lambda}}\rightarrow
   {\D}(\ovl{\Lambda})_{S_0}\rightarrow 
   \Lie \ovl{\Lambda}^\vee\rightarrow 0
 \end{equation*}
 lifts to an exact sequence of formally smooth group schemes
 \begin{equation*}
  0\rightarrow \omega_{\ovl{\Lambda}}\rightarrow
   \E(\ovl{\Lambda})\times_S{S_0}\rightarrow \ovl{\Lambda}^\vee\rightarrow 0
 \end{equation*}
 Via the exponential, it is equivalent to give a local summand $F^1$ of
 ${\D}(\ovl{\Lambda})_S$ which lifts
 $\;\omega_{\ovl{\Lambda}}$, or to give a vector sub-group scheme $V$ of
 $\E(\ovl{\Lambda})$ which lifts $\;\omega_{\ovl{\Lambda}}$. It then
 turns out that $\E(\ovl{\Lambda})/V$ is a 
 $p$-divisible group $\Lambda^\vee$, whose Serre dual is the looked for 
 $p$-divisible group $\Lambda$ over $S$: 
 $\Lambda\times_S S_0=\ovl{\Lambda}$ by duality, and $\omega_{\Lambda}=F^1$.
 

 The formation of $\E(\ovl{\Lambda})$ is functorial (contravariant in 
 $\ovl{\Lambda}$). If we have two $p$-divisible groups $\ovl{\Lambda}, 
 \ovl{\Lambda}'$ and a morphism $\ovl{f}:\;\ovl{\Lambda}=\Lambda\times_S 
 S_0\rightarrow \ovl{\Lambda}'=\Lambda'\times_S S_0$ such that 
 $\D(f)$ respects the $F^1$'s, then the homomorphism 
 $\E(\ovl{f}): \E(\ovl{\Lambda})\rightarrow \E(\ovl{\Lambda}')$ 
 preserves the vector subgroup schemes, hence induces a homomorphism 
 $f: \E(\ovl{\Lambda})/V\cong \Lambda\rightarrow 
 \E(\ovl{\Lambda}')/V'\cong\Lambda'$ which lifts $\ovl{f}$.


 The exponential construction becomes simpler when the ideal of $S_0$ in
 $S$ is of square zero (and endowed with the trivial divided
 powers). This is the case which will be used in the sequel.

\subsection{Crystals and connections.}\label{sub:3.4} 

\begin{para}\label{para:3.4.1}
 We have already encountered the terms ``$F$-crystal'', ``unit-root
 $F$-crystal'' in I.\ref{para:f-cris}, \ref{para-unit-root} in the more
 down-to-earth context of 
 connections. This comes as follows. Assume for simplicity that $S_0$ is
 affine of characteristic $p$ and lifts to a formally smooth $p$-adic
 affine formal scheme $\mathcal{S}$ over $\mathfrak{v}=W$.  Then there
 is an equivalence of categories between finite locally free crystals on
 $(S_0/W)_{\mathrm{cris}}$ and finite locally free
 $\mathcal{O}(\mathcal{S})$-modules with integrable topologically
 nilpotent connection (this interpretation of crystals extends to much
 more general situations, \cf \cite[2.2.2]{jong95:_cryst_dieud}).

 \noindent Using this equivalence of categories, the Katz (covariant) functor
 \begin{equation*}
  \begin{pmatrix}\text{continuous }{\Z}_p\text{-representations} \\
   \text{of } \pi_1^{\mathrm{alg}}(S_0,s_0)
  \end{pmatrix}
  \longrightarrow \text{(unit-root $F$-crystals)}
 \end{equation*}
 is the composite of the following (anti)equivalences of categories
 \begin{align*}
  \begin{pmatrix}\text{continuous } {\Z}_p\text{-representations} \\
   \text{of } \pi_1^{\mathrm{alg}}(S_0,s_0)
  \end{pmatrix}& \longrightarrow
  \text{(ind-etale $p$-divisible groups over $S_0$)}\\
  & \overset{\D}{\longrightarrow}
  \text{(unit-root $F$-crystals over $S_0$)}.
 \end{align*} 
\end{para}

\begin{para}\label{para:3.4.2}
 Let us briefly recall the construction of the connection attached to a
 finite locally free crystal $\mathcal{E}$ on
 $(S_0/W)_{\mathrm{cris}}$. It depends on the technique of formally
 adding divided powers to an ideal (the pd-hull). Let $\Delta_0$ be the
 diagonal in $S_0\times S_0$, and let $\hat\Delta$ be the $p$-completion
 of the pd-hull of the diagonal in $\mathcal{S}\hat\times\mathcal{S}$,
 endowed with the two projections $p_1,p_2$ to $\mathcal{S}$. Then by
 rigidity of crystals, we have for any $n$ an isomorphism
 \begin{equation*}
  p_2^\ast(\mathcal{E}_{S_0,\mathcal{S}/p^n})\;\cong
   \;\mathcal{E}_{\Delta_0,\hat\Delta/p^n}\;\cong
   \;p_1^\ast(\mathcal{E}_{S_0,\mathcal{S}/p^n})
 \end{equation*}
 and at the limit $n\rightarrow \infty$, an isomorphism of finite
 locally free $\mathcal{O}({\hat\Delta})$-modules  
 \begin{equation*}
  \epsilon:\; p_2^\ast(\mathcal{E}_{S_0,\mathcal{S}})\;\cong
   \;p_1^\ast(\mathcal{E}_{S_0,\mathcal{S}})
 \end{equation*}
 satisfying a cocycle condition. Such an isomorphism is the ``Taylor series'' 
 \begin{equation*}
  \epsilon(1\otimes e) = \sum_{n_1,\ldots,n_d\geq
   0}\biggl(\nabla\Bigl(\frac{\partial}{\partial t_1}\Bigr)^{n_1}
   \ldots\nabla\Bigl(\frac{\partial}{\partial t_d}\Bigr)^{n_d}\biggr)(e)
   \otimes \prod_i \frac{\text{``$(1\otimes t_i-t_i\otimes 1)^{n_i}$''}}{n_i!}
 \end{equation*}
 of an integrable connection $\nabla$ on the finite
 locally free $\mathcal{O}(\mathcal{S})$-module
 $\;\mathcal{E}_{S_0,\mathcal{S}}$ (here $t_1,\ldots t_d$ denote local
 coordinates on $\mathcal{S}$).
\end{para}

\begin{para}\label{para:3.4.3}
 In general, the Taylor series converges only in polydiscs of radius
 $| p |^{1/p-1}$, due to the presence of the factorials. One
 says that $\nabla$ is {\it convergent} if its Taylor series converges
 in (open) unit polydiscs. A well-known argument due to Dwork shows that
 $F$-crystals give rise to convergent connections $\nabla$ (\cite{tdd}).

 In order to put the discussion of convergence on proper foundations, we
 need the notion of tube in analytic geometry.
\end{para}


\subsection{Interlude : analytic spaces associated with formal schemes
 and tubes.}\label{sub:3.5}

\begin{para}\label{para:3.5.1} Let us first review the Raynaud-Berthelot
construction of the analytic space attached to a  
$\mathfrak{v}$-formal scheme $\mathfrak{X}$. It will be convenient for
later purpose to formulate the construction in the frame of Berkovich
spaces rather than rigid geometry (we refer to \cite[1.6]{ecfnas} for the
translation).


As in \ref{para:2.2.5}, we assume that $\mathfrak{X}$ is adic, locally
noetherian; moreover, we assume that 
${\mathfrak{X}}_{\mathrm{red}}=\Spec(\mathcal{O}_{\mathfrak{X}}/\mathcal{J})$ is a
separated $k$-scheme locally of finite type (in particular,
$p\in \mathcal{J}$). However, we {\it do not} assume that $\mathfrak{X}$ is
$p$-adic, \ie that the topology of $\mathcal{O}_{\mathfrak{X}}$
is $p$-adic.

 We describe the construction in the affine case ${\mathfrak{X}} =
 \Spf(\mathcal{A})$; the general case follows by gluing. One chooses
 generators $f_1,\ldots, f_r$ of 
 $J=\Gamma({\mathfrak{X}},\mathcal{J})$, and defines, for any $n>0$,
 \begin{equation*}
 \mathcal{B}_n=\mathcal{A}\{T_1,\ldots,T_r\}/(f_1^n-\varpi T_1,\ldots,
 f_r^n-\varpi T_r)
 \end{equation*}
 where $\mathcal{A}\{T_1,\ldots,T_r\}$ stands for the $\varpi$-adic
 completion of $\mathcal{A}[T_1,\ldots,T_r]$. One has
 $\mathcal{B}_n/\varpi\mathcal{B}_n\cong
 (\mathcal{A}/(\varpi,f_1^n,\ldots,f_r^n))[T_1,\ldots,T_r]$, a
 $k$-algebra of finite type; hence $\mathcal{B}_n$ is a
 $\mathfrak{v}$-algebra topologically of finite type (\ie quotient of an
 algebra of the form $\mathcal{A}\{T_1,\ldots,T_{r_n}\}$), so that
 $\mathcal{B}_n[\frac{1}{p}]$ is an affinoid algebra. Let
 $\matheur{M}(\mathcal{B}_n)$ be its Berkovich spectrum. The
 homomorphism $\mathcal{B}_{n+1}\rightarrow \mathcal{B}_n$ sending $T_i$
 to $f_iT_i$ identifies $\matheur{M}(\mathcal{B}_n)$ with an affinoid
 subdomain of $\matheur{M}(\mathcal{B}_{n+1})$ \cite[0.2.6]{crecrasp}.


 The analytic space $\mathfrak{X}^{\mathrm{an}}$ is defined as the union
 $\bigcup_n \matheur{M}(\mathcal{B}_n)$. It does 
 not depend on the choice of $f_1,\ldots, f_r$. This is a {\it
 paracompact (strictly) analytic space} over the fraction 
 field $K$ of $\mathfrak{v}$  (one can replace the nested ``ball-like''
 $\matheur{M}(\mathcal{B}_n)$ by nested ``annulus-like'' 
 spaces in order to obtain locally finite coverings). 
\end{para}

\begin{exs}\label{para:3.5.2}
 (i) if ${\mathfrak{X}}=\Spf(\mathfrak{v}\{t\})$,
 ${\mathfrak{X}}^{\mathrm{an}}$ is the closed unit disc.
 
 (ii) if ${\mathfrak{X}}=\Spf({\mathfrak{v}[[t]]}),
 \;\mathcal{J}=(\varpi, t)$, then $\matheur{M}(\mathcal{B}_n)$ is the 
 closed disc $\matheur{D}_K(0,|\varpi|^{1/n})$, and
 ${\mathfrak{X}}^{\mathrm{an}}$ is the open unit disc. 
\end{exs}

\begin{para}\label{para:3.5.3} Let $S_0$ be a closed subscheme of
$\mathfrak{X}_{\mathrm{red}}$, and let ${\mathfrak{X}}_{S_0}$ be the formal
completion of $\mathfrak{X}$ along $S_0$. The {\it tube of $S_0$ in
$\mathfrak{X}$} is $\mathfrak{X}_{S_0}^{\mathrm{an}}$. It is denoted by
$]S_0[_{\mathfrak{X}}$ (and considered as a subspace of
$\mathfrak{X}^{\mathrm{an}}$). It depends only on $S_{0,\mathrm{red}}$.  

\medskip
 \noindent {\it Example}: if $\mathfrak{X}$ is as in example (i) above
 (so that $\mathfrak{X}_{\mathrm{red}}=\A^1_k$), and $S_0$ is the point $0$,
 then $\mathfrak{X}_{S_0}$ is as in example (ii) above:
 $]0[_{\mathfrak{X}}\;=\matheur{D}(0,1^-)\subset \matheur{D}(0,1^+)$. 

 \noindent More generally, the polydiscs mentioned in \ref{para:3.4.3}
 are the tubes of closed points.
\end{para}


\subsection{Convergent isocrystals.}\label{sub:3.6}     

\begin{para}\label{para:3.6.1} Let $S_0$ be a separated $k$-scheme of finite
 type, and $S_0\hookrightarrow \mathfrak{Y}$ a closed immersion into a
 flat $p$-adic formal $\mathfrak{v}$-scheme that is formally smooth in a
 neighborhood of $S_0$. Let $E$ be a vector bundle on the tube
 $]S_0[_{\mathfrak{Y}}$. An integrable connection $\nabla$ on $E$ is
 {\it convergent} if its Taylor series
 \begin{equation*}
  \epsilon(1\otimes e) =
   \sum_{n_1,\ldots,n_d\geq 0}
   \biggl(\nabla\Bigl(\frac{\partial }{\partial t_1}\Bigr)^{n_1}\ldots
   \nabla\Bigl(\frac{\partial }{\partial t_d}\Bigr)^{n_d}\biggr)(e)
   \otimes \;\prod_i\frac{\text{``$(1\otimes t_i-t_i\otimes 1)^{n_i}$''}}{n_i!}
 \end{equation*}
 is induced by an isomorphism on the tube of the
 diagonal $]S_0[_{\mathfrak{Y}\times \mathfrak{Y}}$
 \begin{equation*}
  \epsilon:\; p_2^\ast(E)\;\cong \;p_1^\ast(E)
 \end{equation*} 
 It amounts to the same to say that for any affine $U$ in $\mathfrak{Y}$
 with local coordinates
 $t_1,\ldots, t_d$, any section $e\in \Gamma(U^{\mathrm{an}},E)$ and any $\eta<1$,
 one has 
 \begin{equation*}
  \Bigl\|\nabla\Bigl(\frac{\partial }{\partial t_1}\Bigr)^{n_1}
   \ldots\nabla\Bigl(\frac{\partial }{\partial t_d}\Bigr)^{n_d}(e)\Bigr\|.
   \eta^{n_1+\ldots + n_d}\rightarrow 0
 \end{equation*}
 where $\|\;\;\|$ denotes a Banach norm on $\Gamma(U^{\mathrm{an}},E)$.

 According to Berthelot, {\it the category of vector bundles
with convergent connection depends only on $S_0$} --- even only on
$S_{0,\mathrm{red}}$ --- (up to canonical equivalence), and is functorial in
$S_0/K$; this is the category of {\it convergent isocrystals on $S_0/K$}
(or, abusively, on $S_0$).
 

Since any separated $k$-scheme locally of finite type can be locally
embedded into a $\mathfrak{X}$ as above, one can define the category of
convergent isocrystals on such a scheme by gluing \cite[2.3.2]{crecrasp}.
\end{para}

\begin{para}\label{para:3.6.2} There is a natural functor from the category
 of $F$-isocrystals on $S_0/\mathfrak{v}$ to the category of 
convergent isocrystals on $S_0/K$ \cite[2.4]{crecrasp}.  The composite of
this functor with the Dieudonn\'e functor is denoted by $\D(-)_\Q$. It
factors through a $\Q_p$-linear functor, still denoted by $\D(-)_\Q:\;$
\begin{equation*}
\Bigl\{\begin{matrix}
	\text{$p$-divisible groups over $S_0$}\\
	\text{up to isogeny}
       \end{matrix}\Bigr\}
 \longrightarrow
 \{\text{convergent isocrystals on $S_0$}\}
\end{equation*}
\end{para}

\begin{para}\label{para:3.6.3} The advantage of convergent isocrystals over
 ``crystals up to isogeny'' lies in their strong 
``continuity property'', which allows to remove the restriction $e<p$ on
ramification (as we shall do henceforth), and work conveniently with
singular schemes $S_0$ (while divided powers present certain
``pathologies'' in the singular case). In our applications, $S_0$ will
be a countable (non-disjoint) union of irreducible projective varieties,
and $\mathfrak{v}$ may be highly ramified.
\end{para}

\begin{para}\label{para:3.6.4} If $S_0$ is the reduced underlying scheme 
 ${\mathfrak{X}}_{\mathrm{red}}$ of a formal scheme $\mathfrak{X}$ as in
 \ref{para:3.5.1}, there 
 is a notion of {\it evaluation of $E$ on $\mathfrak{X}^{\mathrm{an}}$}
 \cite[2.3.2]{crecrasp}: this is a vector bundle on
 ${\mathfrak{X}}^{\mathrm{an}}$ with an integrable connection. 
 

 By functoriality of the construction, it will be enough to deal with the
 affine case ${\mathfrak{X}}=\Spf{\mathcal{A}}$ and to assume that there
 is a closed immersion $S_0=(\Spf{\mathcal{A}})_{\mathrm{red}}\hookrightarrow
 \mathfrak{Y}$ into a formally smooth $p$-adic formal
 $\mathfrak{v}$-scheme. One can construct inductively a compatible family
 of $\mathfrak{v}$-morphisms
 $\Spec(\mathcal{O}_{\mathfrak{X}}/\mathcal{J}^n)\rightarrow \mathfrak{Y}$,
 whence a $\mathfrak{v}$-morphism  $\mathfrak{X} \rightarrow
 \mathfrak{Y}$.  Let $\upsilon :\; ]S_0[_{\mathfrak{X}} 
 ={\mathfrak{X}}^{\mathrm{an}}\rightarrow\; ]S_0[_{\mathfrak{Y}}$ be the
 associated analytic map. The looked for evaluation of $E$ on
 ${\mathfrak{X}}^{\mathrm{an}}$ 
 is the pull-back by $\upsilon$ of the vector bundle with integrable
 connection on $]S_0[_{\mathfrak{Y}}$ defined by $E$.
\end{para}

\begin{para}\label{para:3.6.5} Let now $S$ be a separated scheme in
 $\Nil_\mathfrak{v}$ (it may also be viewed in  $\Nil_W$).
 Let $\mathcal{E}$ be a finite locally free $F$-crystal on $S/W$. This
 provides a vector bundle $\mathcal{E}_S$ on $S$, which 
 only depends on the inverse image $\mathcal{E}\times_W k$ of
 $\mathcal{E}$ on the $k$-scheme $S/p$ (\cf \ref{para:3.2.1}). On the other 
 hand, the convergent isocrystal $E$ on $S/p$ attached to $F$-crystal
 $\mathcal{E}\times_W k$ (\cf \ref{para:3.6.2}) depends only on
 $S_{\mathrm{red}}$.
\end{para}

\begin{para}\label{para:3.6.6} Let us apply this to Dieudonn\'e
 crystals. Let $\mathfrak{X}$ be a formal scheme as in \ref{para:3.5.1},
 and let 
 $\Lambda=\liminj \Lambda_n$ be a $p$-divisible on
 $\mathfrak{X}=\liminj {\mathfrak{X}}_n$. This gives rise to a vector bundle 
 $\displaystyle \D(\Lambda)_{\mathfrak{X}}= \liminj
 \D(\Lambda_n)_{{\mathfrak{X}}_n}$ on $\mathfrak{X}$. On the other 
 hand, we have the vector bundle $\displaystyle
 \D(\Lambda_0)_{\mathfrak{X}}=$ on $\mathfrak{X}$ and a natural
 homomorphism $\D(\Lambda_0)_{\mathfrak{X}}\rightarrow
 \D(\Lambda)_{\mathfrak{X}}$ obtained by pull-back 
 ${\mathfrak{X}}_{\mathrm{red}}={\mathfrak{X}}_0\rightarrow {\mathfrak{X}}$.


When $e\geq p$, this is not an isomorphism in general. However, {\it the
associated morphism of analytic vector bundles
$(\D(\Lambda_0)_{\mathfrak{X}})^{\mathrm{an}}\rightarrow
(\D(\Lambda)_{\mathfrak{X}})^{\mathrm{an}}$ is an isomorphism.  Moreover,
$(\D(\Lambda_0)_{\mathfrak{X}})^{\mathrm{an}}$ is the analytic vector bundle
underlying the evaluation of $\D(\Lambda_0)_\Q$ on
${\mathfrak{X}}^{\mathrm{an}}$}.


Indeed, working locally, and taking in account the very construction of
${\mathfrak{X}}^{\mathrm{an}}$, we reduce to the case when ${\mathfrak{X}}$ is
$p$-adic. Then $\mathfrak{X}/p$ is a $k$-scheme. Using the canonical
divided powers on $(p)$ and the rigidity of crystals,
$\D(\Lambda\times_{\mathfrak{X}}\mathfrak{X}/p)_{\mathfrak{X}}\cong
\D(\Lambda)_{\mathfrak{X}}$. It is thus enough to see that
$\D(\Lambda_0)^{\mathrm{an}}\cong\D
(\Lambda\times_{\mathfrak{X}}\mathfrak{X}/p)_{\mathfrak{X}}^{\mathrm{an}}$. This
follows from the fact that the convergent isocrystals $\D(\Lambda_0)_\Q$
and $\D(\Lambda\times_{\mathfrak{X}}\mathfrak{X}/p)_\Q$ are ``the
same'', since ${\mathfrak{X}}_0=({\mathfrak{X}}/p)_{\mathrm{red}}\;$: indeed, it
is clear that $(\D(\Lambda_0)_{\mathfrak{X}})^{\mathrm{an}}$ (\resp
$\D(\Lambda\times_{\mathfrak{X}}\mathfrak{X}/p)_{\mathfrak{X}}^{\mathrm{an}}$)
coincides with the evaluation of the convergent isocrystal
$\D(\Lambda_0)_\Q$ (\resp
$\D(\Lambda\times_{\mathfrak{X}}\mathfrak{X}/p)_\Q$) on
${\mathfrak{X}}^{\mathrm{an}}$ (\cf also \cite[6.4]{efgonas}, in a slightly less
general context).
\end{para}

\begin{para}\label{para:3.6.7} Let us draw from this two consequences:
 \begin{enumerate}
  \item taking into account the fact that the functor $\D(-)_\Q$ factors
	through the category of $p$-divisible
	groups up to isogeny, we see that if we are given a quasi-isogeny
	$\Lambda_0\rightarrow \Lambda_0'$, then there is a
	canonical isomorphism of vector bundles
	$\bigl(\D(\Lambda'_0)_{\mathfrak{X}}\bigr)^{\mathrm{an}}\rightarrow
	\bigl(\D(\Lambda)_{\mathfrak{X}}\bigr)^{\mathrm{an}}$.
  \item Through the isomorphism
	$\bigl(\D(\Lambda_0)_{\mathfrak{X}}\bigr)^{\mathrm{an}}\rightarrow
	\bigl(\D(\Lambda)_{\mathfrak{X}}\bigr)^{\mathrm{an}}$ (or equivalently,
	through the isomorphism
	$\D\bigl(\Lambda\times_{\mathfrak{X}}
	\mathfrak{X}/p\bigr)_{\mathfrak{X}}^{\mathrm{an}}
	\rightarrow
	\bigl(\D(\Lambda)_{\mathfrak{X}}\bigr)^{\mathrm{an}}$), the 
	analytic vector bundle
	$\bigl(\D(\Lambda)_{\mathfrak{X}}\bigr)^{\mathrm{an}}$ is naturally
	endowed with an integrable connection, which deserves to be
	called the {\it Gauss-Manin} connection attached to $\Lambda$.
 \end{enumerate}

 Indeed, when $\Lambda=\udl{A}[p^\infty]$ for an abelian scheme
 $\udl{A}$ over $\mathfrak{X}$,
 $(\D(\Lambda)_{\mathfrak{X}})^{\mathrm{an}}$ may then be identified
 with the De Rham cohomology bundle
 $\mathrm{H}^1_{\mathrm{DR}}(\udl{A}^{\mathrm{an}}/
 {\mathfrak{X}}^{\mathrm{an}})$
 together with its Gauss-Manin connection.  Here, one can use the
 comparison theorem between crystalline and De Rham cohomology, or the
 theory of universal vectorial extensions and Grothendieck's
 $\natural$-structures, \cf \cite{mazur74:_univer}, \cite{tddc2}.

 In the first alternative, one uses the
 divided powers on $(p)$ and construct a canonical isomorphism of vector
 bundles with connection 
 $\D(\Lambda\times_{\mathfrak{X}}\mathfrak{X}/p)_{\mathfrak{X}}\cong
 \mathrm{H}^1_{\mathrm{DR}}(\udl{A}/{\mathfrak{X}})$.


 In the second alternative, one uses Grothendieck's construction of a
 canonical connection on the universal vectorial extension $\E(\Lambda)$
 of $\Lambda^\vee$ --- \cf \ref{para:3.3.2}; the induced connection on the Lie
 algebra $\D(\Lambda)_{\mathfrak{X}}$ turns out to be the Gauss-Manin
 connection on $ \mathrm{H}^1_{\mathrm{DR}}(\udl{A}/{\mathfrak{X}})$.


 The Gauss-Manin connection is especially significant when
 ${\mathfrak{X}}^{\mathrm{an}}$ is smooth (even though $\mathfrak{X}$
 itself may not be formally smooth).

\newpage
\section{Moduli problems for $p$-divisible groups.}\label{sec:4}

\begin{abst}
 Review of moduli spaces of $p$-divisible groups quasi-isogenous to a
 fixed one modulo $p$, with examples.
\end{abst}
\end{para}


\subsection{Moduli problem for $p$-divisible groups quasi-isogenous to a
 fixed one modulo $p$.}\label{sub:4.1}

\begin{para}\label{para:4.1.1} Here, we work over the discrete valuation
 ring $\mathfrak{v}=W= {\widehat{\Z_p^{\mathrm{ur}}}}$ (the completion of the
maximal unramified extension of $\Z_p$), with uniformizing parameter
$\varpi=p$. Recall that $\Nil_\mathfrak{v}$ denote the category of
locally noetherian $\mathfrak{v}$-schemes $S$ on which $p$ is locally
nilpotent.


Let $\ovl{\Lambda}$ be a fixed $p$-divisible group over
$k={\ovl{\F}}_p$. The following theorem is due to M.  Rapoport and
T. Zink \cite[2.16, 2.32]{psfpg}.
\end{para}

\begin{thm}\label{thm:4.1.2} The functor
\begin{align*}
 \Nil_\mathfrak{v}&\rightarrow \Sets\\
 S&\mapsto \Bigl\{(\Lambda,\rho)\Bigm|
 \begin{array}{l}
  \Lambda \text{: $p$-divisible group on $S$},\\
  \rho\in \qisog(\Lambda\times_S S_{\mathrm{red}},
 \ovl{\Lambda}\times_{{\ovl{\F}}_p}S_{\mathrm{red}})
 \end{array}
 \Bigr\}\Big/\cong
\end{align*}
is representable by a formal scheme $\mathfrak{M}$
over $\mathfrak{v}$; $\mathfrak{M}_{\mathrm{red}}$ is locally of
finite type over ${\ovl{\F}}_p$ and its irreducible components are
projective ${\ovl{\F}}_p$-varieties.
\end{thm}

Of course, the rigidity of $p$-divisible groups up to isogeny (\ref{thm:2.2.3}) is
crucial here. The main difficulty in proving \ref{thm:4.1.2} is to control the
powers of $p$ which are needed in the process of lifting isogenies.


The separated formal scheme $\mathfrak{M}$ is not formally smooth in
general, but has been conjectured to be flat over 
$\mathfrak{v}$.  The group $J$ of self-quasi-isogenies of $\ovl{\Lambda}$
 acts (on the right) on $\mathfrak{M}$ by  $(\Lambda,\rho).j= 
(\Lambda,\rho\circ j$). This group $J$ is actually the group of
 ${\Q}_p$-points of an algebraic group over ${\Q}_p$. 

 We denote by $\mathcal{M}$ the Berkovich analytic space over
 $\C_p$ attached to $\mathfrak{M} : 
 \mathcal{M}={\mathfrak{M}}^{\mathrm{an}}\hat\otimes_K\;\C_p$
 (\cf \ref{sub:3.5}). This is a paracompact (strictly) analytic space. We shall
 see later that it is smooth, and a $p$-adic manifold in the sense of
 I.\ref{para:p-adic-manifold}.

\subsection{Examples.}\label{sub:4.2}
 We take $\ovl{\Lambda} ={\ovl{A}}[p^\infty]$, the $p$-primary
 torsion of an {\it elliptic curve ${\ovl{A}}$ over ${\ovl{\F}}_p$} ($h=2$).

\begin{para}\label{para:4.2.1}
 {\it The ordinary case}, \ie ${\ovl{A}}[p]({\ovl{\F}}_p)\neq 0$. 

 Then $\ovl{\Lambda}\cong\hat {\G}_m \oplus {\Q}_p/{\Z}_p$, and
 $J=({\G_m}_{{\Q}_p})^2$. For any $(\Lambda,\rho)$ in $\mathfrak{M}(S)$,
 the separable rank of $\Ker[p]$ is constant (due to the quasi-isogeny
 $\rho$). Therefore $\Lambda$ is an extension 
 \begin{equation*}
  \hat{\Lambda}\rightarrow \Lambda\rightarrow \Lambda^{\mathrm{et}}.
 \end{equation*}
 The quasi-isogeny $\rho$ respects this extension,
 hence splits into two parts $({\hat\rho},\rho^{\mathrm{et}})$. We have
 ${\hat\Lambda}\cong \hat {\G}_m ,\;\Lambda^{\mathrm{et}}\cong
 {\Q}_p/{\Z}_p$. Up to isomorphism, the pair
 $({\hat\rho},\rho^{\mathrm{et}})$ thus amounts to an element of
 $({\Q}_p^\times/{\Z}_p^\times)^2\cong {\Z}^2$. Therefore
 \begin{equation*}
  \mathfrak{M}=\coprod_{\Z^2}\;\mathfrak{M}^0
 \end{equation*}
 is a disjoint sum of copies of a formal (group) scheme
 $\mathfrak{M}^0$ which parametrizes extensions of ${\Q}_p/{\Z}_p$ by $
 \hat {\G}_m$. It is known that
 \begin{equation*}
  \mathfrak{M}^0 =\hat {\G}_m
 \end{equation*} 
 This identification is given by the
 following recipe. Recall that if $S=\Spec(R),\;R$ local artinian with
 residue field ${\ovl{\F}}_p$ and radical $\mathfrak{m}$
 ($\mathfrak{m}^{n+1}=0$), then $\hat {\G}_m(S)=1+\mathfrak{m}$; in
 particular $\hat {\G}_m(S)$ is killed by $[p^n]$. Let us choose any lift
 $q_n$ of $1/p^n\in {\Q}_p/{\Z}_p$ in $\Lambda$. Then $[p^n]q_n$ is a
 well-defined element $q$ in $ {\hat\Lambda}(S)\cong 1+\mathfrak{m}$,
 which is unchanged if $n$ is replaced by a bigger integer (but another
 choice of the isomorphism ${\hat\Lambda}\cong \hat {\G}_m$ would change
 $q$ into $q^a,\; a\in {\Z}_p^\times$).
 

 \noindent As a formal $W$-scheme, $\mathfrak{M}^0=
 \Spf(\mathfrak{v}[[q-1]])$, with the $(p,q-1)$-adic topology. 
 The associated analytic space $\mathcal{M}^0$ over $\C_p$ is the open
 unit disc $| q-1| <1$.
 
 The parameter $q$ is thus the local modulus for deformations of
 ${\ovl{A}}[p^\infty]$.  
 On the other hand, if
 $p\neq 2$, we have an algebraic local modulus for deformations of
 ${\ovl{A}}$ (say, with level two structure): the Legendre parameter
 $z=\lambda$ around the Legendre parameter $\zeta_{\mathrm{can}}$ of the
 canonical lifting $A_{\mathrm{can}}$ of ${\ovl{A}}$ (\cf 
 I.\ref{para-dwork-para}).
 The $p$-divisible group $A_{\mathrm{can}}[p^\infty]$ splits: its parameter is
 $q=1$. Therefore, in this special case, the Serre-Tate theorem \ref{thm:2.2.4}
 asserts that $q\in
 1+(z-\zeta_{\mathrm{can}}){\widehat{\Z}^{\mathrm{ur}}}[[z-\zeta_{\mathrm{can}}]]$. It has been
 proved by W. Messing and N. Katz that $q$ is the Dwork-Serre-Tate
 parameter discussed in I.\ref{sub-serre-tate-para}, 
 (\cite{q}, \cite{slm}) --- which explains the terminology.
\end{para}

\begin{para}\label{para:4.2.2} {\it The supersingular case}, \ie
 ${\ovl{A}}[p]({\ovl{\F}}_p)= 0$. 
 
 In this case ${\ovl{\Lambda}}
 =\hat{\ovl{\Lambda}}$ is a so-called Lubin-Tate formal group,
 $\;\End\;{\ovl{\Lambda}}=\mathcal{B}_p$, ``the'' maximal order in 
 the (non-split) quaternion algebra $B_p$ over ${\Q}_p$, and $J=B_p^\times$.
 
 
 Here again, there is a discrete invariant: the ``height'' of the
 quasi-isogeny $\rho$, which amounts to the 
 $p$-adic valuation of the norm of its image in $J=B_p^\times$.
 Therefore 
 \begin{equation*}
  \mathfrak{M}=\coprod_{\Z}\;\mathfrak{M}^0
 \end{equation*}
 is a disjoint sum of copies of a formal scheme $\mathfrak{M}^0$
 which parametrizes deformations of the Lubin-Tate group
 $\hat{\ovl{\Lambda}}$. According to Lubin-Tate, $\mathfrak{M}^0\cong
 \Spf(\mathfrak{v}[[t]])$, with the $(p,t)$-adic topology. The
 associated analytic space $\mathcal{M}^0$ over $\C_p$ is the open unit
 disc $|t|<1$.

 Assume for simplicity that ${\ovl{A}}$ is defined over ${\F}_p$. Then
 one can describe the universal formal deformation $\Lambda$ over
 $\Spf{\Z_p[[t]]}$ through its logarithm $\sum b_n(t)X^{p^{2n}}$, which
 is given recursively by
 \begin{equation*}
  b_0(t)=1,\;b_n(t)=\frac{1}{p}\sum_{0\leq m < n}b_m(t)t_{n-m}^{p^{2m}}
 \end{equation*} 

 
The so-called canonical lifting, corresponding to $t=0$, has formal complex
multiplication by $\Z_{p^2}$, \cf \cite{gross94:_equiv_lubin_tate}, \cite{trapmlsasht}.
\end{para}

 
\subsection{Decorated variants.}\label{sub:4.3} 

\begin{para}\label{para:4.3.1} Just as in the complex case
 \ref{para:1.1.4}, it is useful to consider variants or the moduli
 problem \ref{sub:4.1} for ``decorated'' $p$-divisible groups, with prescribed
 endomorphisms and/or polarization.


Instead of
the $\Q$-algebra $B$ of \ref{para:1.1.4}, we shall consider a finite-dimensional semi-simple $\Q_p$-algebra $B_p$, since
the category of $p$-divisible groups up to isogeny is
$\Q_p$-linear. Let $V_p$ be a $B_p$-module
of finite type, and let $G_{\Q_p}$ be the group of 
$B_p$-linear endomorphisms of $V_p$.
\par\noindent  Let $\mathcal{B}_p$ be a maximal order in $B_p$, and let
$L_p$ be a lattice in $V_p$, stable under $\mathcal{B}_p$. 


\end{para}

\begin{para}\label{para:4.3.2} We assume that $\mathcal{B}_p$ acts on our
 $p$-divisible group $\ovl{\Lambda}$. One then redefines $J$ to 
be the group of self-quasi-isogenies of $\ovl{\Lambda}$ which respect
the $\mathcal{B}_p$-action.


For technical reasons, one picks up a lifting $\tilde{\ovl{\Lambda}}$ of
$\ovl{\Lambda}$ with $\mathcal{B}_p$-action to some finite extension of
$W$, and looks at the $\mathcal{B}_p$-module
$F^1_0:=\omega_{\tilde{\ovl{\Lambda}}}$. We denote by $\mathfrak{v}$ the
finite extension $W[\tr(b| F^1_0)]$ of $W$.


One strengthens the moduli problem \ref{sub:4.1} by imposing to our pairs
 $(\Lambda,\rho)$ that $\Lambda$ is endowed with a
 $\mathcal{B}_p$-action $\iota:\mathcal{B}_p\rightarrow \End \Lambda$,
 $\rho$ respects the $\mathcal{B}_p$-action, and moreover a ``Shimura
 type condition'' $\det(\iota(b)| \omega(\Lambda))=\det(b|
 F^1_0)$.

 It follows easily from \ref{thm:4.1.2} \cite[3.25]{psfpg} that this
 moduli problem is {\it representable by a formal scheme, still denoted
 by $\mathfrak{M}$, over $\mathfrak{v}$. It is acted on by
 $J$. $\mathfrak{M}_{\mathrm{red}}$ is locally of finite type over ${\ovl{\F}}_p$
 and its irreducible components are projective
 ${\ovl{\F}}_p$-varieties. }
\end{para}

\begin{para}\label{para:4.3.3} Let us say a word about the polarized variant
 (here, it is safer to assume $p\neq 2$). One assumes 
 that $\mathcal{B}_p$ is endowed with an involution $\ast$, and that
 $L_p$ is autodual for an alternate $\Z_p$-bilinear form such that
 $\an{bv,w}=\an{v,b^\ast w}$. Of course, here, $G_{\Q_p}$ denotes the algebraic
 $\Q_p$-group of $B_p$-linear symplectic similitudes of $V_p$. One
 assumes that $\ovl{\Lambda}$ is endowed with a $\ast$-polarization
 \footnote{a $\ast$-polarization is a symmetric
 $\mathcal{B}_p$-linear quasi-isogeny between the $p$-divisible and its
 Serre dual, endowed with the transposed action of $\mathcal{B}_p$
 twisted by $\ast$}, and modify $J$ to respect the polarization up to a
 multiple in $\Q_p$.
 
 \bigskip
 The corresponding moduli problem is again representable by a formal
 scheme with the same properties. This formal scheme has been conjectured to be flat by Rapoport-Zink; although this has been verified in many cases (\cf \eg \cite{gortz}), there are some counter-examples  (\cf  \cite{pappas}).
\end{para}

\begin{exa}[fake elliptic curves at a critical prime $p$]\label{exa:4.3.4}
 Let us fix a indefinite quaternion algebra $B$
 over $\Q$ and a maximal order $\mathcal{B}$ in $B$ (after Eichler, they
 are all conjugate). A prime $p$ is called {\it critical} if it divides
 the discriminant of $B$, \ie if $B_p=B\otimes_\Q \Q_p$ is the
 (non-split) quaternion algebra over $\Q_p$. Note that $B_p$ then
 contains copies of the unramified quadratic extension on $\Q_p$; we
 denote by $\Q_{p^2}$ one of them. The maximal order (unique up to
 conjugation) can be written $\mathcal{B}_p=\mathcal{B}\otimes
 \Z_p=\Z_{p^2}[\Pi],\;\Pi^2=p,\;\Pi.a=a^\sigma\Pi$ for $a\in
 \Z_{p^2}$. We fix a critical prime $p$.


 We consider abelian surfaces with quaternionic multiplication by
 $\mathcal{B}$. They are sometimes called ``{\it fake elliptic curves}''
 \index{fake elliptic curves@fake elliptic curves}
 because they share many features with elliptic curves, 
 \cf \cite[III, 1]{updcds}.

 We fix a fake elliptic curve $\ovl{A}$ over $\ovl{\F}_p$, and set
 $\Lambda=\ovl{A}[p^\infty]$. This is a special formal
 $\mathcal{B}_p$-modules of height $4$ in the following sense.

 Let $S$ be in $\Nil_\mathfrak{v}$ (or else an inductive limit of such
 schemes --- as a formal scheme), and let $\Lambda$ be a $p$-divisible
 group over $S$. After Drinfeld, one says that $\Lambda$ is a {\it
 special formal $\mathcal{B}_p$-modules of height $4$} if

\begin{enumerate}
 \item $\Lambda$ is infinitesimal ($\Lambda=\hat\Lambda$) of height
       $h=4$,
 \item $\mathcal{B}_p$ acts on $\Lambda$,
 \item via this action, $\omega_\Lambda$ is a locally free
       $(\Z_{p^2}\otimes_{\Z_p}\mathcal{O}_S)$-module of rank $1$.
\end{enumerate}

\noindent We take $V_p=B_p$ where $B_p$ acts by left
multiplication. In this situation, $G_{\Q_p}\cong B_p^\times $. When
$S=\Spec(\ovl{\F}_p)$, there is a unique class of
$\mathcal{B}_p$-isogeny of special formal $\mathcal{B}_p$-modules of
height $4$. The group $J$ of self-quasi-isogenies of a special formal
$\mathcal{B}_p$-module of height $4$ is $J=GL_2(\Q_p)$. This can be
easily read off the Dieudonn\'e module $\otimes\Q$.


In this example, there is again a discrete invariant: the ``height'' of 
the quasi-isogeny $\rho$, which amounts to the 
$p$-adic valuation of the determinant of its image in $J$.
Therefore 
 \begin{equation*}
 \mathfrak{M}=\coprod_{\Z}\;\mathfrak{M}^0
\end{equation*}
 is a disjoint sum of copies of a formal scheme $\mathfrak{M}^0$
 which parametrizes deformations of the special formal module
 $\hat{\ovl{\Lambda}}$. This formal scheme has been described by
 Drinfeld, \cf \ref{para:6.3.4}.
\end{exa}

\begin{para}\label{para:4.3.5} This example and examples \ref{sub:4.2}
 were the sources of the general theory. 

\noindent Let us at once suggest a vague analogy with the
archimedean situation --- especially \ref{para:1.1.4}, \ref{para:1.2.4} ---
 which will become more 
and more precise in the next sections: one can see the $p$-divisible
group $\ovl{\Lambda}$ in characteristic $p$ (\resp $J$) as an analogue
of the infinitesimal oriented real Lie group $\hat U$ (\resp $G(\R)^0$).
    
\newpage
\section{$p$-adic period domains.}\label{sec:5}

\begin{abst}
 Review of $p$-adic symmetric domains and period domains via stability theory.
\end{abst}
\end{para}


\subsection{$p$-adic flag spaces.}\label{sub:5.1}

\begin{para}\label{para:5.1.1} We begin with the $p$-adic analogue of
 the flag space $\mathcal{D}^\vee$ of \ref{para:1.2.3},
 \ref{para:1.2.4}. In \ref{para:1.2.4}, this flag space parametrized the
 lagrangian subspaces $F^1$ in the $\mathrm{H}_{\mathrm{DR}}^1$ of abelian
 varieties of a certain type, marked by a fixed isomorphism
 $\mathrm{H}_{\mathrm{DR}}^1\cong V^\C$.


In the $p$-adic case, $\mathcal{D}^\vee$ will parametrize the Hodge
filtration $F^1$ of Dieudonn\'e modules of $p$-divisible groups of a
certain type (\ref{para:3.3.1}).
\end{para}

\begin{para}\label{para:5.1.2} More precisely, let us consider a moduli
 problem for $p$-divisible groups as in \ref{sub:4.1}, possibly
 decorated as in \ref{sub:4.3} (from which we follow the notation). The
 decoration reflects on the 
 Dieudonn\'e module $\D(\ovl{\Lambda})_\Q=\D(\ovl{\Lambda})\otimes \Q$:
 one has an isomorphism of $B_p$-module (with symplectic structure, in
 the polarized case):
 \begin{equation*}
  \D(\ovl{\Lambda})_\Q\cong
   V_p\otimes_{\Q_p}\widehat{\Q_p^{\mathrm{ur}}}.
 \end{equation*}
 We fix such an isomorphism, and consider that
 $G_{\Q_p}\otimes_{\Q_p}\widehat{\Q_p^{\mathrm{ur}}}$ acts on
 $\D(\ovl{\Lambda})_\Q$. But $\D(\ovl{\Lambda})_\Q$ has an extra
 structure: the $\sigma$-linear Frobenius automorphism. The group $J$ is
 nothing but the {\it subgroup} of $G_{\Q_p}(\widehat{\Q_p^{\mathrm{ur}}})$ of
 elements commuting with Frobenius.
\end{para}

\begin{para}\label{para:5.1.3} Let $\mathfrak{v}$ be the finite extension of
 $W={\Z_p^{\mathrm{ur}}}$ considered in \ref{para:4.3.2} and let $K$ be its 
 fraction field. We denote here by $F^1_0\subset V_p^{\C_p}=V_p\otimes
 \C_p$ the $\C_p$-span of the space
 $\omega_{\tilde{\ovl{\Lambda}}}\subset V_p\otimes_{\Q_p}K$ introduced in
 \ref{para:4.3.2}. We set $G_{\C_p}=G_{\Q_p}\otimes_{\Q_p}\C_p$.

 \noindent The flag variety $\mathcal{D}^\vee$ of all
 $G_{\C_p}$-conjugates of $F^1_0$ in $V_p^{\C_p}=\D(\ovl{\Lambda})
 \otimes_W\C_p$ is a homogeneous space 
 \begin{equation*}
  \mathcal{D}^\vee = P\backslash G_{\C_p}
 \end{equation*}
 for a suitable parabolic subgroup $P$. There is an
 ample line bundle $\mathcal{L}$ on $\mathcal{D}^\vee$
 which is homogeneous under the derived group $G_{\C_p}^{\mathrm{der}}$;
 \index{000GCpder@$G_{\C_p}^{\mathrm{der}}$}
 we denote by $\mathcal{D}^\vee \subset \P(W)$ an 
 associated $G_{\C_p}^{\mathrm{der}}$-equivariant projective embedding (if
 $G_{\C_p}=GL_{2g},\;\mathcal{D}^{\vee} = \Grass(g,2g)$, one 
 can just take the $SL_{2g}$-equivariant Pl\"ucker embedding).

 Note that, there is an obvious right action of $J$ on
 $\mathcal{D}^{\vee}$, which factors through the adjoint group
 $J^{\mathrm{ad}}:=\Im\bigl(J\rightarrow G_{\C_p}^{\mathrm{ad}}\bigr)$. 
\end{para}

\begin{para}\label{para:5.1.4} In our three basic examples (\ref{para:4.2.1}, \ref{para:4.2.2}, \ref{exa:4.3.4}), we
have $G_{\C_p}=GL_2$, hence $\mathcal{D}^\vee=\P^1_{\C_p}$.    
\end{para}


\subsection{Symmetric domains via stability theory.}\label{sub:5.2}


\begin{para}\label{para:5.2.1} In the $p$-adic case, one cannot define symmetric spaces by a positivity condition as in \ref{para:1.2.3}, \ref{para:1.2.4}.
One uses instead the Hilbert-Mumford notion of semi-stability, following an idea of M. van der Put and H. Voskuil
\cite{ssatsagoalf}. 


\end{para}

\begin{para}\label{para:5.2.2}
 We set $J'=J\cap G_{\C_p}^{\mathrm{der}}$. This is
 again the group of $\Q_p$-points of an algebraic group over $\Q_p$. Let
 $T\cong ({\G_m}_{\Q_p})^j$ be a maximal $\Q_p$-split torus in $J'$. The
 {\it semi-stable locus} of $\mathcal{D}^\vee$ with respect to $T$-action
 is the open subset of points $x\in \mathcal{D}^\vee$ such that there is
 a $T$-invariant function of $\P(W)$ which does not vanish on $x$.


 One can check whether $x$ is semi-stable with respect to $T$ using the
 Hilbert criterion: let $\lambda$ be any cocharacter (= one-parameter
 subgroup) of $T$. By properness of $\mathcal{D}^\vee$, the map $z\in
 \G_m\mapsto x.\lambda(z)\in \mathcal{D}^\vee$ extends to $\A^1$, \ie the
 limit $\displaystyle x_0=\lim_{z\rightarrow 0} x.\lambda(z)$ exists. It
 is clearly fixed by $\G_m$, thus the fiber $\mathcal{L}_{x_0 }$
 corresponds to a character $z\mapsto z^{-\mu(x,\lambda)}$ of
 $\G_m$. Then $x$ is semi-stable if and only if the Mumford invariant
 $\mu(x,\lambda)$ is $\geq 0$ for every $\lambda$.
\end{para}

\begin{para}\label{para:5.2.3} One defines the {\it symmetric space
 $\mathcal{D}$ to be the intersection of the semi-stable loci 
 of $\mathcal{D}^\vee$ with respect to all maximal $\Q_p$-split tori $T$
 in $J'$.} This is a $p$-adic manifold in the sense of
 I.\ref{para:strict-analytic}.  It is
 also called the {\it $p$-adic period domain} (associated with the moduli
 problem).


 Since $J$ acts on the set of tori $T$ by conjugation, it is clear that
 $\mathcal{D}$ is stable under the $J$-action on $\mathcal{D}^{\vee}$.
\end{para}

 
\subsection{Examples.}\label{sub:5.3}

\begin{para}\label{para:5.3.1} In our three basic examples, we have
$G_{\C_p}=GL_2,\;\mathcal{D}^\vee=\P^1$.


In example \ref{para:4.2.1}, $J'\subset J=(\G_m)^2$ is the one-dimensional torus
$T$ with elements $(z^{-1},z)$. Its action on a point $x$ is by
$z^2$-scaling; the limit point $x_0$ is always $0$, except if
$x=\infty$. The only non-semi-stable point is $\infty$, hence
$\mathcal{D}=\A^1$.



In example \ref{para:4.2.2}, $J'=SL_1(B_p)\subset J=B_p^\times$. There is no
 non-trivial $\Q_p$-split torus in $J'$, hence
 $\mathcal{D}=\mathcal{D}^{\vee}=\P^1$.


\end{para}

\begin{para}\label{para:5.3.2} We now turn to example \ref{exa:4.3.4}. In this
 example, $J=GL_2(\Q_p)$ and $J'=SL_2(\Q_p)$. All maximal 
 $\Q_p$-split tori are conjugate $g^{-1}Tg$ of the above $T$. Therefore,
 $\mathcal{D}=\P^1\setminus \bigcup \infty.g$ 
 is the {\it Drinfeld space} over $\C_p$ 
 \begin{equation*}
  \mathcal{D}=\Omega_{\C_p} :=\P^1_{\C_p}\setminus \P^1(\Q_p)
 \end{equation*} 

 To be consistent, we consider the right action of
 $PGL_2$-action on $\P^1$, which is deduced from the customary left 
 action by the rule $x.g=g^{-1}.x$. 

 The Drinfeld space is actually defined over $\Q_p$, \ie is a
 $\Q_p$-analytic manifold $\Omega$. 
 It is closely related to the {\it Bruhat-Tits tree}
 $\mathcal{T}$ of $J^{\mathrm{ad}}=PGL_2(\Q_p)$. The vertices of this
 tree correspond to closed discs in 
 $\Q_p$; two such vertices are connected by an edge if for the
 corresponding discs $\matheur{D}',\matheur{D}''$, one has
 $\matheur{D}'\subset \matheur{D}''$ and the  
 radius of $\matheur{D}''$ is $p$ times the radius of $\matheur{D}'$
 (each vertex has $p+1$ 
 neighbors). The group $PGL_2(\Q_p)$ acts on $\mathcal{T}$ (on 
 the right, say).
 \vspace*{5mm}
 %%% PICTURE
 \begin{figure}[h]
%  \setlength{\unitlength}{1pt}
  \begin{picture}(120,120)(0,0)
   \put(0,0){\includegraphics[scale=0.5,keepaspectratio,clip]{fig2-5a.eps}}
%   \put(100,50){\circle*{3}}
%   \put(100,50){\line(0,1){20}}\put(100,70){\circle*{3}}
%   \put(100,50){\line(2,-1){20}}\put(120,40){\circle*{3}}
%   \put(100,50){\line(-2,-1){20}}\put(80,40){\circle*{3}}
%   %
%   \put(100,70){\line(2,1){12}}\put(112,76){\circle*{3}}
%   \put(100,70){\line(-2,1){12}}\put(88,76){\circle*{3}}
%   \put(120,40){\line(0,-1){12}}\put(120,28){\circle*{3}}
%   \put(120,40){\line(2,1){12}}\put(132,46){\circle*{3}}
%   \put(80,40){\line(0,-1){12}}\put(80,28){\circle*{3}}
%   \put(80,40){\line(-2,1){12}}\put(68,46){\circle*{3}}
%   %
%   \multiput(112,76)(0,4){3}{\circle*{.3}}
%   \multiput(112,76)(4,-2){3}{\circle*{.3}}
%   \multiput(88,76)(0,4){3}{\circle*{.3}}
%   \multiput(88,76)(-4,-2){3}{\circle*{.3}}
%   \multiput(120,28)(4,-2){3}{\circle*{.3}}
%   \multiput(120,28)(-4,-2){3}{\circle*{.3}}
%   \multiput(132,46)(0,4){3}{\circle*{.3}}
%   \multiput(132,46)(4,-2){3}{\circle*{.3}}
%   \multiput(80,28)(4,-2){3}{\circle*{.3}}
%   \multiput(80,28)(-4,-2){3}{\circle*{.3}}
%   \multiput(68,46)(0,4){3}{\circle*{.3}}
%   \multiput(68,46)(-4,-2){3}{\circle*{.3}}
  \end{picture}\vspace{-1ex}
%  \setlength{\unitlength}{1.5pt}
  \caption{$\mathcal{T}$ for $p=2$}
  \label{fig-bruhat-tits}
 \end{figure}

 For any disc $\matheur{D}$ centered in $\Q_p$, let $\eta_\matheur{D}$
 denote its Berkovich 
 ``generic point'' (which is a point in Berkovich's affine line $\A^1$
 over $\Q_p$). Then the geometric realization of 
 $\mathcal{T}$ may be identified with the closed subset 
 \begin{equation*}
  |\mathcal{T}|= \{\eta_{\matheur{D}(a,r^+)} \mid a\in \Q_p,\;0<r<\infty\}
  \,\subset \,\Omega
 \end{equation*}
 (a $PGL_2(\Q_p)$-equivariant embedding), and the set of vertices with
 the subset 
\begin{equation*}
 \mathrm{ver}\, \mathcal{T}=
  \{\eta_{\matheur{D}(a,r^+)} \mid a\in \Q_p,\;r\in p^\Z\}\, \subset
  \,|\mathcal{T}| 
\end{equation*}
 According to Berkovich, $|\mathcal{T}|$ is a retract of $\Omega$
 \cite{staagonf}; in particular, $\Omega$ is simply connected.

 On the other hand, $\Omega$ is the generic fiber of a $p$-adic formal
 scheme $\hat\Omega$ whose reduction modulo $p$ is an 
 infinite tree of projective lines, with dual graph $\mathcal{T}$ (\cf
 \cite[I]{updcds}).


 The Drinfeld space $\Omega$ is a $p$-adic analogue of the Poincar\'e
 double half-plane $\P^1_{\C}\setminus \P^1(\R)$. Note that the latter
 space could also be defined using stability, in the same way (with
 $J'=SL_2(\R)$).  Another similarity: according to \cite{staagonf},
 there is a $PGL_2(\Q_p)$-invariant metric on $\Omega$ (which 
 extends the standard metric on $|\mathcal{T}|$). 
\end{para}


\newpage
\section{The $p$-adic period mapping and the Gauss-Manin connection.}
\label{sec:6}

\begin{abst}
 Rapoport-Zink's period mapping as classifying map for the Hodge
 filtration in the Dieudonn\'e module. Basic properties: equivariance,
 \'etaleness. Its relation to the Gauss-Manin connection. Examples and
 formulas.
\end{abst}

\subsection{Construction.}\label{sub:6.1} 

\begin{para}\label{para:6.1.1} The construction of the $p$-adic period
 mapping is parallel to the complex case: one 
 attaches to a point $s\in \mathcal{M}$ representing a pair $(\Lambda,
 \rho)$ the point $\mathcal{P}(s)$ of $\mathcal{D}^\vee$ 
which parametrizes the notch $F^1$ of the Hodge filtration of the
 Dieudonn\'e module of $\Lambda$. The idea of the construction is
 already present in Grothendieck's talk
 \cite{grothendieck70:_group_barsot_tate}.
\end{para}

\begin{para}\label{para:6.1.2} More precisely, notation being as in
 \ref{para:4.1.1}--\ref{para:4.3.1}, let us consider the universal object 
 $(\udl{\Lambda},\udl{\rho})$ on the moduli formal scheme
 $\mathfrak{M}$, and the Dieudonn\'e crystal $\D(\udl{\Lambda})$. Its
 evaluation $\D(\udl{\Lambda})_{\mathfrak{M}}$ on $\mathfrak{M}$ is a
 vector bundle on $\mathfrak{M}$.
\end{para}

\begin{lem}\label{lem:6.1.3} The quasi-isogeny $\udl{\rho}$ induces a
 trivialization of the associated analytic vector bundle on
$\mathcal{M}={\mathfrak{M}}^{\mathrm{an}}_{\C_p} $:
\begin{equation*}
 (\D(\udl{\Lambda})_{\mathfrak{M}})^{\mathrm{an}}\cong
 V_p^{\C_p}\otimes_{\C_p}\mathcal{O}_\mathcal{M}
\end{equation*}


\end{lem}

\begin{proof}
 $\rho$ is a quasi-isogeny between $\udl{\Lambda}
 \times_{\mathfrak{M}}{\mathfrak{M}}_{\mathrm{red}}$ and
 $\ovl{\Lambda}\times_{\ovl{\F}_p}{\mathfrak{M}}_{\mathrm{red}}$. According to
 \ref{para:3.6.7} (i), by the rigidity of convergent isocrystals, it induces a
 canonical isomorphism $(\D(\udl{\Lambda})_{\mathfrak{M}})^{\mathrm{an}}\cong
 \bigl(\D(\ovl{\Lambda}\times_{\ovl{\F}_p}
 {\mathfrak{M}}_{\mathrm{red}})_{\mathfrak{M}}\bigr)^{\mathrm{an}}.$
 On the other hand, since the formation of $\D(\Lambda)$ commutes with
 base change (\ref{para:3.2.2}), we have
 \begin{equation*}
 \bigl(\D(\ovl{\Lambda}\times_{\ovl{\F}_p}
 {\mathfrak{M}}_{\mathrm{red}})_{\mathfrak{M}}\bigr)^{\mathrm{an}} =
 \D(\ovl{\Lambda})\otimes_W 
 \mathcal{O}_\mathcal{M}=V_p\otimes_{\Q_p}\mathcal{O}_\mathcal{M}.
 \end{equation*}
\end{proof}

\begin{para}\label{para:6.1.4}
 On the other hand, we have a locally direct summand
 $F^1=(\omega_{\udl{\Lambda}})^{\mathrm{an}}$ in
 $(\D(\udl{\Lambda})_{\mathfrak{M}})^{\mathrm{an}}\cong
 V_p\otimes_{\Q_p}\mathcal{O}_\mathcal{M}$ (\ref{para:3.3.1}). This
 determines an 
 $\mathcal{M}$-point of the flag space $\mathcal{D}^\vee$, \ie an
 analytic mapping 
 \begin{equation*}
  \mathcal{P}:\mathcal{M}\;\longrightarrow \;\mathcal{D}^\vee
 \end{equation*}

 This is called the {\it period mapping} for the moduli problem
 (decorated or not). It is clear from the construction that $\mathcal{P}$ is
 $J$-{\it equivariant}.


 If $\mathcal{M}^0$ is any connected component of $\mathcal{M}$, it is
 stable under $J'\subset J$, hence the induced period mapping
 $\mathcal{P}:\;\mathcal{M}^0\;\longrightarrow \;\mathcal{D}^\vee\;$ is
 $J'$-equivariant (the $J'$-action factors through ${J'}^{\mathrm{ad}}$).
\end{para}

\subsection{Properties.}\label{sub:6.2}  

\begin{thm}\label{thm:6.2.1} The period mapping $\mathcal{P}$ is etale. Its image lies in the period domain $\mathcal{D}$. 


\end{thm}
The first assertion reduces, via an infinitesimal criterion of
etaleness, to the ``essential surjectivity part'' of the
Grothendieck-Messing theorem (\ref{thm:3.3.3}), in the case of an ideal
of square zero, \cf \cite[5.17]{psfpg}.  It corresponds to what is called
the ``local Torelli property'' in the complex case.


The second assertion is a
consequence of a theorem of B. Totaro on the equivalence of semi-stability and ``weak admissibility'' \cite{tpipht}.
The second assertion may be viewed as a
$p$-adic analogue of the ``Riemann relations'' (\ref{para:1.1.3}).

\begin{cor}\label{cor:6.2.2} $\mathcal{M}$ is smooth, and any component
 $\mathcal{M}^0$ is a connected $p$-adic manifold in the sense of
 I.\ref{para:p-adic-manifold}.
\end{cor}
\begin{para}\label{para:6.2.3}
 It is very likely that $\mathcal{P}$, viewed as a mapping
 $\mathcal{M}\rightarrow \mathcal{D}$, is surjective.  This is closely
 related to Fontaine's conjecture on the realizability of filtered
 Dieudonn\'e modules by $p$-divisible groups. Although recent work by
 C. Breuil \cite{breuil99:_repres}, \cite{breuil00:_group} settles the
 Fontaine conjecture at least for 
 finite residue fields of characteristic $\neq 2$ (and in other cases),
 it is not clear to the author to which extent the surjectivity
 conjecture is now established.


 On the other hand, the fibers of the period mapping are relatively
 well-understood, \cf \cite[5.37]{psfpg}. 
\end{para}

\begin{para}\label{para:6.2.4}
 The argument in the proof of \ref{lem:6.1.3} shows a little more. Recall from
 \ref{para:3.6.7} (ii) that the analytic vector bundle
 $(\D(\udl{\Lambda})_{\mathfrak{M}})^{\mathrm{an}}$ on $\mathcal{M}$ comes
 together with an integrable connection $\nabla_{\mathrm{GM}}$ (the Gauss-Manin
 connection).
\end{para}

 \begin{lem}\label{lem:6.2.5} The trivialization 
  \begin{equation*}
   (\D(\udl{\Lambda})_{\mathfrak{M}})^{\mathrm{an}}\cong
   V_p^{\C_p}\otimes_{\C_p}\mathcal{O}_\mathcal{M}
  \end{equation*}
  of \ref{lem:6.1.3} induces a trivialization of the Gauss-Manin connection:  
  \begin{equation*}
   \bigl(
    (\D(\udl{\Lambda})_{\mathfrak{M}})^{\mathrm{an}}
    \bigr)^{\nabla_{\mathrm{GM}}} 
    \cong V_p^{\C_p}.
  \end{equation*}

  The point is that the convergent isocrystal
  $\D(\Lambda\times_{\mathfrak{M}}{\mathfrak{M}}_{\mathrm{red}})$ is constant, 
  since it comes by base change, from $\D(\ovl{\Lambda})_\Q$.
 \end{lem}
 This useful lemma will allow us to give explicit expression of the
 period mapping in terms of quotients of solutions of differential
 equations, as in the complex case. 


\subsection{Examples and formulas.}\label{sub:6.3} 

\begin{para}\label{para:6.3.1} We start with example \ref{para:4.1.1}
 ($p$-divisible groups of elliptic curves with ordinary reduction).
 The extension
 \begin{equation*}
  1\rightarrow {\hat\G_m}\rightarrow \Lambda\rightarrow \Q_p/\Z_p\rightarrow 1
 \end{equation*}
 parametrized by $q\in 1+pW$ gives rise, by contravariance, to an
 extension of Dieudonn\'e modules (\cf \ref{para:3.1.1}) 
 \begin{equation*}
  0\rightarrow W\rightarrow \D(\Lambda)\rightarrow W(-1)\rightarrow 0
 \end{equation*}
 which splits canonically due to the rigidity of Dieudonn\'e crystals:
 $\D(\Lambda)=\D(\ovl{\Lambda})=W\oplus W(-1)$.  We denote by
 $e_0,e_{-1}$ the corresponding canonical basis of $\D(\Lambda)$. The
 image $\tau=\mathcal{P}(q)\in W[\frac{1}{p}]$ of $q$ under the period
 mapping is, by definition, the unique element such that $\tau
 e_0+e_{-1}$ generates $F^1$. It turns out that
 \begin{equation*}
  \tau=\log q,
 \end{equation*} 
 thus
 \begin{equation*}
  \mathcal{M}^0=\matheur{D}(1,1^-)
  \xrightarrow{\mathcal{P}=\log} \A^1 =\mathcal{D}.
 \end{equation*}

 Let us indicate a sketch of proof of this remarkable formula (the
 formula goes back to Messing-Katz \cite{q}, \cite{slm}).  The argument
 is more transparent if 
 we think in terms of one-motives instead of elliptic curves: that is,
 rather than viewing $\Lambda$ as the $p$-divisible attached to some
 elliptic curve over $W$, one remarks that it is also the $p$-divisible
 group attached to the one-motive $M=[\Z
 \xrightarrow{1\mapsto q} \G_m]$. Then $\D(\Lambda)$ can be identified
 with $H_{\mathrm{DR}}(M)$, the dual of the Lie algebra of $G^\natural$, where
 $[\Z \rightarrow G^\natural]$ denotes the universal vectorial extension
 of $M$ (\cf \cite{tdh3}); this extension splits canonically
 ($G^\natural=\G_a\times \G_m$), and this splitting agrees with the
 above splitting of $\D(\Lambda)$. We are thus reduced to prove the
 following purely algebraic fact: let us view $M=[\Z
 \xrightarrow{1\mapsto q} \G_m]$ as a
 one-motive over $K((q-1))$; then in terms of the canonical basis
 $e_0,e_{-1}$ of $H_{\mathrm{DR}}(M)$, $F^1\subset
 H_{\mathrm{DR}}(M)$ is generated by the 
 vector $(\log q)e_0+e_{-1}$. This well-known fact is easily proved by
 transcendental means, replacing $K$ by $\C$.
\end{para}

\begin{para}\label{para:6.3.2} Since the basis $e_0,e_{-1}$ is horizontal
 under the Gauss-Manin connection $\nabla$ on $\matheur{D}(1,1^-)\;$ 
 (\ref{lem:6.2.5})\footnote{which is induced by the usual Gauss-Manin
 connection attached to the Legendre elliptic pencil, \cf
 \ref{para:3.6.6} (ii).}, $\tau$ is a quotient 
 of solutions of the Gauss-Manin connection. In terms of the Legendre
 parameter $z=\lambda$ (around the parameter $\zeta_{\mathrm{can}}$ of the
 canonical lifting $A_{\mathrm{can}}$ of the fixed ordinary elliptic curve
 $\ovl{A}$), the equality $\tau = \log q$ means that $q\in
 1+(z-\zeta_{\mathrm{can}})W[[z-\zeta_{\mathrm{can}}]]$ is the
 Dwork-Serre-Tate parameter 
 discussed in I.\ref{sub-serre-tate-para} (for $p\neq 2$).


 This allows to give explicit formulas for $\mathcal{P}$ as a function
 of the algebraic parameter $z$ rather than $q$. For instance, let us
 assume that $p\equiv 1\;(4)$. The disc $| z-\frac{1}{2}| <1$ is
 an ordinary disc. The canonical lifting is $\zeta_{\mathrm{can}}=\frac{1}{2}$
 (complex multiplication by $\Z[\sqrt{-1}]$).  The Gauss-Manin
 connection is represented by the hypergeometric differential operator
 with parameter $(\frac{1}{2},\frac{1}{2},1)$.  Let us recall the
 setting, which we have already met in I.\ref{para-supersingular-disk}.

 We consider the symplectic basis $\omega_1=\bigl[\frac{dx}{2y}\bigr],\;
 \omega_2=\bigl[\frac{(2x-1)dx}{4y}\bigr]$ of relative De Rham
 cohomology.  It has 
 the virtue to be, at $z=\frac{1}{2}$, a basis of eigenvectors for the
 complex multiplication. The Gauss-Manin connection satisfies
 \begin{equation*}
  \omega_2=2z(z-1)\nabla\Bigl(\frac{d}{dz}\Bigr)\omega_1 +
  \frac{4z-5}{6}\;\omega_1.
 \end{equation*} 
 The first row of the
 solution matrix $Y$ normalized by $Y\bigl(\frac{1}{2}\bigr)=\id$ is given by
 \begin{align*}
  y_{11} &=
  F\Bigl(\frac{1}{4},\frac{1}{4},\frac{1}{2};(1-2z)^2\Bigr) +
  \frac{1}{2}(1-2z)F\Bigl(\frac{3}{4},\frac{3}{4},\frac{3}{2};(1-2z)^2\Bigr),\\
  y_{12} &=
  (1-2z)F\Bigl(\frac{3}{4},\frac{3}{4},\frac{3}{2};(1-2z)^2\Bigr).
 \end{align*}
 We have to relate ${\omega_1}_{|_{1/2}},\;
 {\omega_2}_{|_{1/2}}$ to the symplectic basis $e_0,e_{-1}$. Because
 $z=1/2$ gives the canonical lifting, the canonical splitting of the
 Dieudonn\'e module is given by the complex multiplication: $e_0,e_{-1}$
 is a basis of eigenvectors. In fact, taking into account the fact that
 we deal with symplectic bases, we find
 \begin{equation*}
  {\omega_1}_{|_{1/2}}=\frac{\Theta}{2}e_{-1},\quad
   {\omega_2}_{|_{1/2}}=\frac{2}{\Theta}e_0,
 \end{equation*}
 where $\Theta\in W^\times$ denotes the Tate constant
 (for $z=\frac{1}{2}$) discussed in I.\ref{para-tate-para},
 \ref{sub-computation}. (the ambiguity --- by a 
 factor in $\Z_p^\times$ --- is removed since the basis $e_0,e_{-1}$ has 
 been fixed).  The period mapping $\tau=\tau(z)$ is given by the quotient
 of the entries in the first row of the matrix
 $Y.\begin{pmatrix}
     {\Theta/2}&0\\
     0&{2/\Theta}
    \end{pmatrix}$, that is to say:
 \begin{equation*}
  \tau(z)=\Bigl(\frac{\Theta^2}{4}\Bigr)\frac{(1-2z)F(\frac{3}{4},\frac{3}{4},
   \frac{3}{2};(1-2z)^2)}{F(\frac{1}{4},\frac{1}{4},
   \frac{1}{2};(1-2z)^2)+\frac{1}{2}(1-2z)F(\frac{3}{4},
   \frac{3}{4},\frac{3}{2};(1-2z)^2)}
 \end{equation*}
 
 Note that since the target of the period mapping is $\A^1$, the
 denominator does not vanish on $\matheur{D}(\frac{1}{2},1^-)$. Note
 also that, in 
 analytic terms, the etaleness of $\mathcal{P}$ (\ref{thm:6.2.1}) reflects the
 fact that $\frac{d\tau}{dz}$ does not vanish, being the quotient of
 linearly independent solutions of a linear differential equation of
 rank two.
\end{para}

\begin{para}\label{para:6.3.3}
 We now turn to example \ref{thm:4.1.2} ($p$-divisible groups of
 elliptic curves with supersingular reduction).  In this case, the period
 mapping 
 \begin{equation*}
  \mathcal{M}^0\cong \matheur{D}(0,1^-) \xrightarrow{\mathcal{P}}
  \P^1=\mathcal{D} 
 \end{equation*}
 was thoroughly investigated by Gross-Hopkins
 \cite{gross94:_equiv_lubin_tate}, \cite{trapmlsasht}. 
 In the special case considered in \ref{thm:4.1.2}. ii)
 (the fixed supersingular elliptic curve $\ovl{A}$ being defined over
 $\ovl{\F}_p$), they obtained the following formula
 \begin{equation*}
  \mathcal{P}(t)= \lim_{n\rightarrow \infty}
  p^{n}b_{2n}(t)/p^{n+1}b_{2n+1}(t) \in \Q_p[[t]]
 \end{equation*}
 and a closed formula was proposed by J. K. Yu in \cite[11]{Yu95}.

 \noindent In terms of the algebraic Legendre parameter
 $z=\lambda$ in a supersingular disc, it is still true, via
 \ref{lem:6.2.5}, that 
 $\tau=\mathcal{P}(z)$ is a quotient of solutions of the Gauss-Manin
 connection. This allows to give explicit formulas for $\mathcal{P}$ as a
 function of the algebraic parameter $z$. For instance, let us assume
 that $p\equiv 3\;(4)$. The disc $| z-\frac{1}{2}| <1$ is a
 supersingular disc. In terms of a symplectic basis of eigenvectors for
 the action of $[\sqrt{-1}]$ on $V_p^{\C_p}$, we get the following
 formula for the Gross-Hopkins period mapping
 \begin{equation*}
  \tau(z)=\kappa\;\frac{(1-2z)
   F\bigl(\frac{3}{4},\frac{3}{4},\frac{3}{2};(1-2z)^2\bigr)}
   {F\bigl(\frac{1}{4},\frac{1}{4},\frac{1}{2};(1-2z)^2\bigr)
   +\frac{1}{2}(1-2z)F\bigl(\frac{3}{4},\frac{3}{4},\frac{3}{2};(1-2z)^2\bigr)}
 \end{equation*}
 where $\kappa$ is a constant depending on the choice of the basis (a
 natural choice would be to take this basis inside the $p$-adic Betti
 lattice discussed in I.\ref{sub-p-adic-betti-lattice-supersing}; the
 corresponding constant $\kappa$ is then the one occuring in the proof
 of I.\thmref{thm-cm-period}.


 Note that since the target of the period mapping is $\P^1$, the
 denominator vanishes somewhere on $\matheur{D}(\frac{1}{2},1^-)$. 
 \par\noindent Note also that the period mapping in the complex
 situation\footnote{for the Legendre elliptic
 pencil, and computed with help of a basis of cycles --- with coefficients
 in the Gauss integers --- which are eigenvectors for complex
 multiplication at z=1/2.} is given, up to a constant factor, by
 {\it the same hypergeometric formula, viewed as a complex analytic function}. 
\end{para}

\begin{para}\label{para:6.3.4}
 At last, let us turn to example \ref{exa:4.3.4} ($p$-divisible groups of fake
 elliptic curves at a critical prime $p$). In this case, the period mapping 
 \begin{equation*}
  \mathcal{M}^0 \xrightarrow{\mathcal{P}}\Omega_{\C_p} =\mathcal{D}
 \end{equation*}
 is an {\it isomorphism}. This is a reinterpretation
 of Drinfeld's work, \cf \cite{psfpg} (and \cite[II8]{updcds}). Actually,
 Drinfeld establishes an isomorphism at the level of formal schemes 
 \begin{equation*}
  {\mathfrak{M}}^0 \cong \hat{\Omega}\,
  \hat\otimes_{\Z_p}\widehat{\Z_p^{\mathrm{ur}}}
 \end{equation*}
 which is $J^1$-equivariant (where $J^1$ denotes the
 image in $J^{\mathrm{ad}}=PGL_2(\Q_p)$ of the elements of $J$ whose
 determinant is a $p$-adic unit).
\end{para}

\begin{para}\label{para:6.3.5} Recall that the $p$-adic manifold
 $\Omega_{\C_p}$ is simply connected (in the usual topological
 sense). However, the above isomorphism and the modular property of
 $\mathfrak{M}$, allowed Drinfeld to construct a tower of finite etale
 connected Galois coverings of $\Omega_{\C_p}$ (or even of
 $\Omega_{\widehat{\Q_p^{\mathrm{ur}}}}$), as follows.


 One considers the universal special formal module $\udl{\Lambda}$ of
 height $4$ over
 $\mathfrak{M}^0\cong\;\hat\Omega\;\hat\otimes\widehat{\Z_p^{\mathrm{ur}}}$. For
 any $n\geq 1$, $\Ker[p^n]$ is a finite locally free formal group scheme
 of rank $p^{4n}$ over
 $\hat\Omega\;\hat\otimes\widehat{\Z_p^{\mathrm{ur}}}$. This gives rise to a
 finite etale covering of the analytic space
 $\Omega_{\widehat{\Q_p^{\mathrm{ur}}}}$ (fibered in rank-one
 $\mathcal{B}_p/p^n\mathcal{B}_p$-modules).


 To get a connected Galois covering, one looks at the $\Pi$-action
 (recall from \ref{exa:4.3.4} that $\mathcal{B}_p=\Z_{p^2}[\Pi],\;\Pi^2=p$): one
 has $\Ker[p^n]=\Ker[\Pi^{2n}]$, and the complement of $\Ker[\Pi^{2n-1}]$
 in $\Ker[\Pi^{2n}]$ provides a finite etale Galois covering $\Sigma^n$
 of $\Omega_{\widehat{\Q_p^{\mathrm{ur}}}}$ with group
 $(\mathcal{B}_p/p^n\mathcal{B}_p)^\times$. When $n$ grows, these form a
 projective system of Galois coverings (with Galois group the profinite
 completion of $\mathcal{B}_p^\times$). This tower is equivariant with
 respect to the $GL_2(\Q_p)$-action.

 Over $\C_p$, $\Sigma^n$ decomposes into finitely many copies of
 a connected $p$-adic manifold $\Upsilon^n$ which is a finite etale
 Galois covering of $\Omega_{\C_p}$ with group
 $SL_1(\mathcal{B}_p/p^n\mathcal{B}_p)$.  This can be seen either by
 using various maximal commutative $\Q_p$-subalgebras of $B_p$ and
 reasoning as in \cite[1.4.6]{varshavsky98:_shimur}, or by global means, using the
 \v{C}erednik-Drinfeld theorem (\ref{para:7.4.7}), or else by the
 properties of the so-called determinant map \cite[II]{tpuosc}.


Up to now, there is no other method for constructing non-abelian finite
etale coverings of $\Omega_{\C_p}$. The coverings $\Upsilon^n$ remain
rather mysterious: are they simply-connected? Is any connected finite
etale covering of $\Omega_{\C_p}$ a quotient of some $\Upsilon^n$? Does
there exist an infinite etale Galois covering of analytic manifolds
$\Upsilon\rightarrow \Omega_{\C_p}$ with Galois group
$SL_1(\mathcal{B}_p)$?


Drinfeld's construction works for any $GL_n$ over any local field $L$,
replacing $\Omega$ by the complement ${}_L\Omega^{(n-1)}$ of
$L$-rational hyperplanes in $\P^{n-1}$.


\subsection{De Jong's viewpoint on $p$-adic period mappings.}\label{sub:6.4}

 J. de Jong \cite{efgonas} has proposed an interpretation of the
 period mapping which clarifies a lot the nature of its fibers. We
 briefly present this interpretation under the following mild simplifying
 assumption:

\begin{itemize}
 \item[$\ast$] that there is a family of 
 affinoid domains $(X_i)$ of $\mathcal{M}^0$ such that the
 $\mathcal{P}(X_i)$ (which are finite unions of affinoid domains in
 $\mathcal{D}$ since
 $\mathcal{P}$ is etale) form an admissible covering of $\mathcal{D}$. 
\end{itemize}

This assumption is fulfilled in our three basic examples (but not
always, \cf \cite[5.53]{psfpg}).

 Under $(\ast)$, De Jong shows (\loccit\footnote{of course
 $p=\ell$ throughout \S\S  6 and 7 of \loccit }, intr. and \ref{sub:7.2}) that
 there is an etale local system of 
 $\Q_p$-spaces $\mathcal{V}$ on $\mathcal{D}$ such that
 $\mathcal{M}^0$ is a component of the space of
 $\Z_p$-lattices in $\mathcal{V}$. Here, an ``etale local system
 of $\Q_p$-spaces'' is given by the data
 \begin{equation*}
  \mathcal{V}=(\{U_i\hookrightarrow \mathcal{D}\}, \mathcal{V}_i, \phi_{ij})
 \end{equation*}
 where $U_i\hookrightarrow \mathcal{D}$ are open immersions,
 $\mathcal{V}_i=\underset{n}{\limproj}
 \;\mathcal{V}_i/p^n\mathcal{V}_i$ are etale local systems of
 $\Z_p$-lattices (each $\mathcal{V}_i/p^n\mathcal{V}_i$ being a finite
 locally free sheaf of $\Z/p^n\Z$-modules on the etale site of $U_i$),
 and $ \phi_{ij}: (\mathcal{V}_i\otimes \Q_p)_{|
 U_i\times_\mathcal{D}U_j} \rightarrow (\mathcal{V}_j\otimes \Q_p)_{|
 U_i\times_\mathcal{D}U_j} $ are isomorphisms satisfying the usual
 cocycle condition (\cf \cite[4]{efgonas}).


 Of course, in the ``decorated case'', $\mathcal{V}$ inherits the
 relevant decoration.

 It follows from this interpretation that in example \ref{thm:4.1.2}, the
 fibers of the period mapping are in natural bijection with the set of
 vertices of the Bruhat-Tits tree $\mathcal{T}$.

 Similarly, in example \ref{exa:4.3.4}, one can see that the fibers of
 $\mathcal{P}$ are in bijection with the 
 one-point set $SL_1(B_p)/SL_1(\mathcal{B}_p)$, \cf \ref{para:7.4.4}.
\end{para}

\newpage
\section{$p$-adic uniformization of Shimura varieties.}\label{sec:7}

\begin{abst}
 Application of the theory of $p$-adic period mappings: Rapoport-Zink's
 $p$-adic uniformization theorem. The case of global uniformization.
 \v{C}erednik-Drinfeld uniformization of Shimura curves and
 other examples.
\end{abst}

\subsection{$p$-integral models of Shimura varieties.}\label{sub:7.1}

\begin{para}\label{para:7.1.1} Let $\Sh$ be a Shimura variety of PEL type as
 in \ref{para:1.1.4}. We recall the setting of \ref{para:1.1.4}: $B$ is
 a simple finite-dimensional $\Q$-algebra, with a positive involution
 $\ast$, $V$ is a $B$-module 
 of finite type, endowed with an alternate $\Q$-bilinear form such that
 $\an{bv,w}=\an{v,b^\ast w}$; 
 $G$ denotes the $\Q$-group of
 $B$-linear symplectic similitudes of $V$; $\mathcal{B}$ is a maximal
 order in $B$ stable under $\ast$, and 
 $L$ is a lattice in $V$, stable under $\mathcal{B}$ and autodual for
 $\an{\;,\;}$.

 The (non-connected) Shimura variety $\Sh$ is defined over the reflex
 field $E$ (which is a number field in $\C$).  There are variants of
 $\Sh$ in which the principal congruence subgroup of level $N$ of
 $G(\hat\Z)$ is replaced by any open compact subgroup $C$ of $G(\A_f)$,
 where $\A_f=\hat\Z\otimes \Q$ is the ring of finite adeles of $\Q$. The
 complex points of $\Sh$ admit the following adelic description
\begin{equation*}
\Sh(\C)=C\backslash \bigl((G(\R)^{\mathrm{ad}}.h_0) \times G(\A_f)\bigr) /G(\Q)
\end{equation*}
(\cf \ref{para:1.2.4} about the one-parameter group $h_0$). 
\end{para}

\begin{para}\label{para:7.1.2} We set $B_p=B\otimes_{\Q}
 \Q_p,\;\mathcal{B}_p=\mathcal{B}\otimes_\Z \Z_p,$ and so on. We denote
 by $C_p$ the stabilizer of $L_p$ in $G(\Q_p)$

 {\it We assume that $C$ is of the form $C^pC_p$} for some open compact
 subgroup $C^p$ of $G(\A_f^p)$, where $\A_f^p=(\prod_{\ell\neq
 p} \Z_\ell)\otimes \Q$. 


 In the case of a congruence subgroup of level $N$, this condition means
 that {\it the level is prime to $p$}. 

 We fix an embedding $v: E\hookrightarrow
 \C_p$, denote by $E_v$ the $v$-completion of $E$, by $\mathfrak{v}$ the
 discrete valuation ring
 $\widehat{\Z^{\mathrm{ur}}}.\mathcal{O}_{E_v}=
 \widehat{\Z^{\mathrm{ur}}}.\mathcal{O}_E$ 
 (compositum in $\C_p$), and by $K$ the fraction field of $\mathfrak{v}$.
\end{para}

\begin{para}\label{para:7.1.3} Let $S$ be a scheme. The objects of the the
 category of $S$-abelian schemes {\it up to prime to $p$ 
isogeny} are the abelian schemes over $S$; morphisms between $A $ and
$A'$ are global sections of $ \udl{\Hom}_S(A, A')\otimes_{\Z}\Z_{(p)}$,
where $\Z_{(p)}$ is the ring of rational numbers with denominator prime
to $p$. This is a $\Z_{(p)}$-linear category. Its relation to
$p$-divisible groups is given by the following elementary lemma:
\end{para}

\begin{lem}\label{lem:7.1.4} The faithful functor $A\mapsto A[p^\infty]$
\begin{equation*}
 \{\text{$S$-abelian schemes up to prime to $p$
  isogeny}\}\rightarrow \{\text{$p$-divisible groups over $S$}\}
\end{equation*}
is conservative, \ie reflects isomorphisms.
 


\end{lem}
The following result (which we state in a rather vague manner) is due to
R. Kottwitz \cite{possvoff} (Compare with \ref{para:1.2.4}.).

\begin{pro}\label{pro:7.1.5}
 For sufficiently small $C^p$, there is a moduli scheme $\Sh$ over
 $\mathcal{O}_{E_v}$ for isomorphism classes of abelian 
 varieties $A$ up to prime to $p$ isogeny, with $\mathcal{B}$-action 
 (with Shimura type condition), together with a $\Q$-homogeneous
 principal $\ast$-polarization and a class of $B$-linear symplectic
 similitudes $\mathrm{H}^1_{\mathrm{et}}(A,\A_f^p)\rightarrow \A_f^p\otimes V$
 modulo the group $C^p$. The generic fiber $\Sh
 \otimes_{\mathcal{O}_{E_v}}E_v$ is a finite sum of copies of
 $\Sh\otimes_E {E_v}$.
\end{pro}

In the case of a congruence subgroup of level $N$, this condition that
$C^p$ is sufficiently small means that $N$ is sufficiently large (and
kept prime to $p$).

\subsection{The uniformization theorem.}\label{sub:7.2}

\begin{para}\label{para:7.2.1}
 One fixes a point of $\Sh(\ovl{\F}_p)$, hence a
 $\ovl{\F}_p$-abelian variety $\ovl{A}$ up to prime to $p$ isogeny with
 $\mathcal{B}$-action, principal $\ast$-polarization, and level
 structure. We consider the $p$-divisible group
 $\ovl{\Lambda}={\ovl{A}}[p^\infty]$ (with $\mathcal{B}_p$-action,
 principal 
 $\ast$-polarization, level structure). We denote by $\ovl{G}$ the
 $\Q$-group of self-quasi-isogenies of ${\ovl{A}}$ respecting the
 additional structure. It is clear that $\ovl{G}(\Q_p)\subset J$; on the
 other hand $\ovl{G}(\A_f^p)\subset G(\A_f^p)$.

 We also consider an auxiliary lifting $\tilde{\ovl{A}}$ of $\ovl{A}$ (with
 $\mathcal{B}$-action) over some finite extension $\mathcal{O}_{E'}$ of
 $\mathcal{O}_E$ in $\mathfrak{v}$, and set
 $\tilde{\ovl{\Lambda}}=\tilde{\ovl{A}}_\mathfrak{v}[p^\infty]$.  We can
 then identify $\D(\ovl{\Lambda})\otimes_W K$ with
 $\mathrm{H}^1_{\mathrm{DR}}(\tilde{\ovl{A}})\otimes_{\mathcal{O}_{E'}}
 K$, and the submodule $\omega_{\tilde{\ovl{\Lambda}}}$ with
 $F^1\mathrm{H}^1_{\mathrm{DR}}(\tilde{\ovl{A}})\otimes_{\mathcal{O}_{E'}}
 K$.  Thus if the complex lagrangian space 
 $F^1_0$ of \ref{para:1.1.4} is taken to be
 $F^1\mathrm{H}^1_{\mathrm{DR}}(\tilde{\ovl{A}})\otimes_{\mathcal{O}_{E'}}
 \C$ (for some complex embedding of $E'$ 
 extending the natural complex embedding of $E$), the complex and
 $p$-adic Shimura type conditions ``agree'' (\ref{para:1.1.4}, \ref{para:4.3.2}). We can
 consider the corresponding decorated moduli problem for $p$-divisible
 groups, and the formal moduli scheme $\mathfrak{M}$ over $\mathfrak{v}$
 as in \ref{sub:4.3}.
 

 \noindent Moreover, the flags spaces $\mathcal{D}^\vee$ considered in 
 the complex and $p$-adic situations agree, \ie come from the same flag 
 space defined over the reflex field $E$.


\end{para}

\begin{para}\label{para:7.2.2}
 One defines a morphism of functors on $\Nil_\mathfrak{v}$: 
 \begin{equation*}
  {\mathfrak{M}} \rightarrow \Sh
 \end{equation*}
 as follows \cite[6.14]{psfpg}.
 Let $S$ be in $\Nil_\mathfrak{v}$ and let us consider a pair
 $(\Lambda,\rho)\in \mathfrak{M}(S)$. By rigidity of $p$-divisible
 groups up to isogeny (\ref{thm:2.2.3}), the
 $\mathcal{B}_p$-quasi-isogeny $\rho 
 \in \qisog(\Lambda\times_S S_{\mathrm{red}},
 \ovl{\Lambda}\times_{{\ovl{\F}}_p}S_{\mathrm{red}})$ lifts to a unique 
 $\mathcal{B}_p$-quasi-isogeny $\tilde{\rho}\in \qisog(\Lambda,
 \tilde{\ovl{\Lambda}}\times_\mathfrak{v}S)$. This can be algebraized in
 a weak sense: there is an $S$-abelian scheme $A_{(\Lambda,\rho)}$ up to
 prime to $p$ isogeny, with $\mathcal{B}$-action and with $p$-divisible
 group $\Lambda $, and a quasi-isogeny $A_{(\Lambda,\rho)}\rightarrow
 \tilde{\ovl{A}}\times_{\mathcal{O}_{E'}}S$ which induces $\tilde{\rho}$
 at the level of $p$-divisible groups. Moreover $A_{(\Lambda,\rho)}$ is
 unique in the category of $S$-abelian schemes up to prime to $p$ isogeny
 with $\mathcal{B}$-action (\ref{lem:7.1.4}), and its formation is
 functorial in 
 $(\Lambda,\rho)$. One then defines the required morphism of functors on
 $\mathfrak{M}(S)$ by setting $(\Lambda,\rho)\mapsto A_{(\Lambda,\rho)}$
 (endowed with the polarization and level structure inherited from $\ovl{A}$).

 This defines a morphism of formal schemes, hence a morphism
 of associated analytic spaces over $\C_p$:
 \begin{equation*}
  \mathcal{Q}:\;\mathcal{M}\;\rightarrow\; \Sh_{\C_p}^{\mathrm{an}}.
 \end{equation*}
\end{para}

\begin{para}\label{para:7.2.3}
 It turns out that the decorated abelian varieties (up to prime to $p$
 isogeny) which are {\it isogenous} to $\ovl{A}$ form in a natural way a
 {\it countable union $Z=\bigcup Z_n$ of projective irreducible
 subvarieties $Z_n$ of $\;\Sh_{\ovl{\F}_p}$}, each of them
 meeting only finitely many members of the family (\cite[6.23, 6.34]{psfpg}).
 One can slightly generalize the notion of tubes (\ref{sub:3.5}) and
 define the tube $]Z[\subset \Sh_{\C_p}^{\mathrm{an}}$. We denote by $S$ the
 connected component of $]Z[$ which contains the modular point $s$ of our
 chosen lifting $\tilde{\ovl{A}}$.
\end{para}

\begin{thm}\label{thm:7.2.4}
 For sufficiently small $C^p$, there is a connected component
 $\mathcal{M}^0$ of $\mathcal{M}$ such that the restriction
 $\mathcal{Q}:\;\mathcal{M}^0\rightarrow S$ is a topological covering of
 $S$. More precisely, $S$ is the (right) quotient of $\mathcal{M}^0$ by
 some torsion-free discrete subgroup of $J$.
\end{thm}

\cf \cite[6.23, 6.31]{psfpg}. Roughly speaking, the proof of \ref{thm:7.2.4}
involves two steps: using the Serre-Tate theorem \ref{thm:2.2.4}, one
shows that 
the morphism of $\mathfrak{v}$-formal schemes $\mathfrak{M}$
$\rightarrow \widehat{\Sh}$ defined above is
formally etale. One concludes by a careful study of this morphism on
geometric points.

In our examples \ref{sub:4.2}, $S$ is an open unit
disc, and the map $\mathcal{Q}$ is an isomorphism (in \ref{para:4.2.1}, it is
induced by $q\mapsto z=\lambda$).

\begin{para}\label{para:7.2.5} Let us consider the universal decorated
 abelian scheme $\udl{A}$ over $\Sh$, and the vector bundle
 $\mathcal{H}=\mathrm{H}^1_{\mathrm{DR}}(\udl{A}/\Sh)$ with its
 Gauss-Manin connection 
 $\nabla_{\mathrm{GM}}:\mathcal{H}\rightarrow \Omega^1_{\Sh}\otimes
 \mathcal{H}$. 
\end{para}

\begin{cor}\label{cor:7.2.6}  Viewed as a $p$-adic connection and
 restricted to $S$, the Gauss-Manin connection 
 ${\nabla_{\mathrm{GM}}}_{| S}$ comes from a representation of
 $\pi_1^{\mathrm{top}}(S,s)$ on $V^{\C_p}\;$ (hence it satisfies
 Cauchy's theorem, \cf I.\ref{sub-connection}).
\end{cor}

\begin{proof}
 Let $\tilde{S}$ be the universal covering of $S$. According to
 I.\ref{sub-connection}), we have to show that the pull-back of 
 $\nabla_{\mathrm{GM}}$ on $\tilde{S}$ is trivial ($\cong V^{\C_p}\otimes
 \mathcal{O}_{\tilde{S}})$. By \ref{thm:7.2.4} and \ref{lem:6.2.5}, this
 already holds on the quotient $\mathcal{M}^0$ of $\tilde{S}$.
\end{proof}

 \noindent {\it Remark.} Using the Weil descent data of \cite[6.21]{psfpg},
 one can show that the representation is defined over a
 finite unramified extension $E'_v$ of $E_v$, \ie comes from a
 representation of $\pi_1^{\mathrm{top}}(S,s)$ on $V\otimes E'_v$.

\begin{para}\label{para:7.2.7}
 The assumption $C=C^pC_p$ is necessary to define the morphism
 $\mathcal{Q}: \mathcal{M}^0\rightarrow S$. The assumption ``$C^p$ small
 enough'' is much less important. Theorem \ref{thm:7.2.4}. holds without it,
 except that the discrete subgroup of $J$ is then no longer torsion-free
 [the point is that there is always a normal subgroup ${C'}^p$ of finite
 index which is small enough, one applies the theorem to ${C'}^p$, and one
 passes to the quotient by ${C'}^p/C^p$]. However, to give a modular
 interpretation of the (ramified) quotient of $\mathcal{M}^0$ in this
 case, one needs the notion of $p$-adic orbifold (to be discussed in
 III); moreover the Gauss-Manin connection may acquire logarithmic
 singularities along the branched locus.

\subsection{Global uniformization}\label{sub:7.3}

We say that there is {\it global uniformization}
in case $S$ is (the analytification of) a whole connected component of
the Shimura variety $\Sh$. This occurs when the reduction mod $p$ of
the decorated abelian varieties $A$ parametrized by $\Sh$ form a 
{\it single isogeny class}. 

In the situation of global uniformization, the group $\ovl{G}$ of
self-quasi-isogenies of ${\ovl{A}}$ (with decoration) is an
{\it inner form of $G$}.

In view of our analogy \ref{para:4.3.5}, this situation can be expected to be
closest to the complex case (in the notation of \ref{para:1.2.4}, the $\hat A$'s
are all isomorphic to $\hat U$). Indeed, we then get a commutative
diagram of $p$-adic manifolds (similar to \ref{para:1.2.4}):
\begin{equation*}
 \begin{CD}
  \mathcal{M}^0@>\mathcal{P}>> \mathcal{D} @. \subset\;\mathcal{D}^{\vee}\\
  @V\mathcal{Q}VV @VV\mathcal{Q}V @.\\
  S@>\mathcal{P}>> \mathcal{D}/\Gamma.
 \end{CD}
\end{equation*}
 where the horizontal maps are etale, $\mathcal{Q}$ denotes the quotient
 maps, and $\Gamma$ is an arithmetic subgroup in the semi-simple group
 $\ovl{G}^{\mathrm{ad}}$. In all known examples, $\mathcal{P}$ is
 actually an isomorphism and $\mathcal{M}^0=\tilde S$ 
 \cite[6]{psfpg}, \cite{tpuosc}, \cite{varshavsky98:_shimur},
 \cite{varshavsky98:_shimur_ii}.


Here, the archimedean---non archimedean correspondence goes beyond a mere
analogy: indeed, here and in \ref{para:1.2.4}, what is denoted by $S$ (\resp
$\mathcal{D}^\vee $) is the --- complex or $p$-adic --- analytification of
the same algebraic variety.

\medskip
\noindent {\it Remark.} To avoid any ambiguity, let us point out
that we consider here only phenomena of global uniformization which can
be interpreted in terms of moduli spaces of $p$-divisible groups. This
does not account for all cases of global $p$-adic uniformization of
Shimura varieties. For instance, it is known that the elliptic modular
curve $X_0(p)$ is a Mumford curve over $\Q_p$, whereas the $p$-adic
uniformization of \ref{thm:7.2.4} (or variants with level $p$ structure)
is local.
\end{para}



\subsection{Example: the \v{C}erednik-Drinfeld uniformization of
Shimura curves at a critical prime.}\label{sub:7.4}

\begin{para}\label{para:7.4.1}
 The first historical example of global uniformization in the above
 sense is the \v{C}erednik uniformization of ``Shimura
 curves'' (modular curves for fake elliptic curves) at a critical prime
 $p$, \cf \cite{uoacbdsopwcq}.  \v{C}erednik's method was
 of group-theoretic nature and did not involve formal groups. A modular
 proof was subsequently proposed by Drinfeld \cite{copsr} (\cf also
 \cite[III]{updcds}). For examples, and related topics, we refer to
 \cite{van89:_les_shimur}, \cite{van92:_discr_mumfor}, \cite{van92:_unifor}.
\end{para}

\begin{para}\label{para:7.4.2}
 We fix an indefinite quaternion algebra $B$ over $\Q$, a maximal order
 $\mathcal{B}$ in $B$, and a positive involution $\ast$ of
 $\mathcal{B}$. We consider the standard representation $V=B$ of $B$, so
 that $G=GL_B(V)\cong B^\times$.

 We refer to the foundational paper \cite{shimura67:_const} for the
 properties of the (connected) Shimura curve ${\mathcal{X}^+(N)}$
 of level $N$ attached to $B$. It parametrizes principally 
 $\ast$-polarized abelian surfaces with multiplication by $\mathcal{B}$
 (fake elliptic curves, \cf \ref{exa:4.3.4}), with level $N$ structure.
 We recall the following properties of ${\mathcal{X}^+(N)}$, which 
 actually characterize it (\loccit 3.2):

\begin{enumerate}
\renewcommand{\theenumi}{\alph{enumi}}
 \item ${\mathcal{X}^+(N)}$ is a projective geometrically connected smooth
       curve defined over the cyclotomic field $\Q(\zeta_N)$,
 \item ${\mathcal{X}^+(N)}(\C)= \mathfrak{h}/\Gamma^+(N)$, that is to
       say: ${\mathcal{X}^+(N)}$ is the quotient of the Poincar\'e upper
       half plane by the image $\Gamma^+(N)=PSL_1(\mathcal{B})(N)$ of
       the group $\{g\in (1+N\mathcal{B})^\times, \Nr(g)>0\}$ in
       $PSL_2(\R)$, where $\Nr$ stands for the reduced norm,
 \item for any imaginary quadratic field $K$ in $B$ such that
       $\mathcal{O}_K\subset \mathcal{B}$, let $z$ be its fixed point in
       $\mathfrak{h}$. Then the extension of $\Q(\zeta_N)$ generated by the
       image of $z$ in ${\mathcal{X}^+(N)}$ is the class field over $K$ with
       conductor $N\mathcal{B}\cap \mathcal{O}_K$.
\end{enumerate}

 The commutative diagram of \ref{para:1.2.4}, involving the complex period 
 mapping, specializes to the following 

 \begin{equation*}
  \begin{CD}
   \tilde{S} @>\mathcal{P}>> \mathcal{D}=\mathfrak{h} @.
   \subset \; \mathcal{D}^{\vee}=\P^1\\
   @V\mathcal{Q}VV @VV\mathcal{Q}V\\
   S={\mathcal{X}^+(N)}_\C^{\mathrm{an}}@>\mathcal{P}>>\mathcal{D}/\Gamma^+(N)
  \end{CD}
 \end{equation*}
 where the horizontal maps are isomorphisms (and the
 vertical maps are topological coverings if $\Gamma^+_p(N)$ is
 torsion-free\footnote{if this group is not torsion-free $\tilde S$ 
 is the universal covering in the sense of orbifolds.}).

 There is a useful variant $\mathcal{X}^\ast$ of
 $\mathcal{X}^+={\mathcal{X}^+(1)}$, already considered by Shimura
 (\cite[3.13]{shimura67:_const}, in which 
 $\Gamma^+$ is replaced by the image $\Gamma^\ast$ in $PSL_2(\R)$ of the
 group $\{g\in B^\times \mid g\mathcal{B}=\mathcal{B}g,\,\Nr(g)>0\}$:
 \begin{equation*}
  \mathcal{X}^\ast(\C)= \mathfrak{h}/\Gamma^\ast.
 \end{equation*}
 It is defined over $\Q$. More generally, there is a
 Shimura curve $\mathcal{X}_\Gamma$ attached to any 
 congruence subgroup $\Gamma$ in $B^{\times^{+}}/\Q^\times$, where
 $B^{\times^{+}}$ denotes the group of elements of positive 
 reduced norm.
\end{para}

\begin{para}\label{para:7.4.3}
 It is known that ${\mathcal{X}^+(N)}$ has
 good reduction at any prime $p$ which does not divide the discriminant 
 $d(B)$ of $B$ nor the level $N$ (Y. Morita). If instead $p$ is critical
 (but still {\it does not divide $N$}), as we shall 
 assume henceforth, \v{C}erednik's theorem shows that
 ${\mathcal{X}^+(N)}_{\C_p}$ is a {\it Mumford curve} uniformized by 
 the Drinfeld space $\Omega_{\C_p}=\P^1_{\C_p}\setminus \P^1(\Q_p)$ (at
 least if $N$ is large enough); in fact, $\Sh$ becomes a 
 Mumford curve over some finite unramified extension of $\Q_p$ (note that
 $\Q(\zeta_N)\subset {\Q_p^{\mathrm{ur}}}$).  

 To be more explicit, let $\ovl{B}$ be the definite quaternion
 $\Q$-algebra which is ramified at the same primes as $B$ except $p$:
 $\ovl{B}_p\cong M_2(\Q_p)$. Let $\ovl{\mathcal{B}}$ be a maximal order in
 $\ovl{B}$. Let $\Gamma^+_p(N)$ be the image in $PGL_2(\Q_p)$ of the
 group $\{g\in (1+N\ovl{\mathcal{B}}[\frac{1}{p}])^\times \mid \ord_p(\Nr(g))
 \;{\rm even}\}$. Let $\Gamma^\ast_p$ be the image in $PGL_2(\Q_p)$ of
 the normalizer $\{g\in {\ovl{B}}^\times \mid g\ovl{\mathcal{B}}[\frac{1}{p}] =
 \ovl{\mathcal{B}}[\frac{1}{p}]g\}$.
 These are discrete subgroups, and one has the {\it \v{C}erednik
 uniformization}:
 \begin{align*}
  {\mathcal{X}^+(N)}_{\C_p}^{\mathrm{an}}&\cong\Omega_{\C_p}/\Gamma^+_p(N),\\
  {\mathcal{X}^{\ast}}_{\C_p}^{\mathrm{an}}&\cong\Omega_{\C_p}/\Gamma^\ast_p.
 \end{align*} 
 For recent applications to $p$-adic $L$-functions, Heegner points,
 etc... \cf \cite{bertolini98:_heegn_l_cered}.
\end{para}

\begin{para}\label{para:7.4.4}
 The case of $\mathcal{X}^{\ast}$ is the one originally considered in 
 \cite{uoacbdsopwcq}, except that \v{C}erednik uses the
 adelic language. In order to make the translation, it is useful to keep
 in mind the following few facts (\cf \cite[pp. 40, 99 and passim]{adadq}).


 For critical $p$, the unique maximal order
 $\mathcal{B}_p$ of $B_p$ is $\{b\in B_p \mid \Nr(b)\in
 \Z_p\}$. Hence, we have an equality $SL_1(\mathcal{B}_p) = SL_1(B_p)$
 and a chain of compact normal subgroups
 \begin{equation*}
  PSL_1(\mathcal{B}_p) = PSL_1(B_p)\subset PGL_1(\mathcal{B}_p) \subset
   PGL_1(B_p)
 \end{equation*}
 with $PGL_1(B_p)/PSL_1(\mathcal{B}_p)\cong \Q_p^\times/(\Q_p^\times)^2$
 (which is of type $(2,2)$ if $p\neq 2$, of type $(2,2,2)$ if $p=2$). 

 On the other hand, the normalizer of $\mathcal{B}_\ell$ in
 $B_{\ell}^{\times}$ is $B_{\ell}^{\times}$ if $\ell$ is critical, and
 $\mathcal{B}_\ell^\times.\Q_{\ell}^{\times}$ otherwise. It follows that  

 \begin{align*}
  \Gamma^\ast&=\{g\in B^{\times^+}\mid \forall \ell
  \nmid d(B), \; g\in \mathcal{B}_\ell^\times.\Q_\ell^\times\}/\Q^\times\\
  \Gamma^\ast_p&=\{g\in \ovl{B}^\times\mid \forall
   \ell \nmid d(B), \; g\in
   \ovl{\mathcal{B}}_\ell^\times.\Q_\ell^\times\}/\Q^{\times}\\
  \Gamma^+(N)&= \{g\in B^{\times^+}\mid \forall \ell, \; g\in
   (1+N\mathcal{B}_\ell)^\times.\Q_\ell^{\times}\}/\Q^{\times}\\
  \Gamma^+_p(N)&= \Bigl\{g\in \ovl{B}^\times\Bigm|
   \begin{array}{l}
    \ord_p(\Nr(g))\;{\rm even};\\
    \forall \ell \neq p,\; g\in
    (1+N\ovl{\mathcal{B}}_{\ell})^\times.\Q_\ell^\times\\ 
   \end{array}
  \Bigr\}\Big/\Q^{\times}
 \end{align*}
 If $N$ is prime to $d(B)$, this can also be rewritten
 \begin{align*}
 \Gamma^+(N)&= \Bigl\{g\in B^{\times^+}\Bigm|
  \begin{array}{l}
   \forall\ell\mid d(B),\, \ord_\ell(\Nr(g))\;{\rm even};\\
   \forall\ell\nmid d(B),\, g\in (1+N\mathcal{B}_\ell)^\times.\Q_\ell^\times\\
  \end{array}
  \Bigr\}\Big/\Q^\times \\
  \Gamma^+_p(N)&= \Bigl\{g\in \ovl{B}^\times\Bigm|
  \begin{array}{l}
   \forall\ell\mid d(B),\, \ord_\ell(\Nr(g))\;{\rm even};\\
   \forall\ell\nmid d(B),\,
    g\in (1+N\ovl{\mathcal{B}}_{\ell})^{\times}.\Q_{\ell}^{\times}\\
  \end{array}
  \Bigr\}\Big/\Q^{\times}
 \end{align*}
 In particular, for $N=1$, we see that there are canonical bijections 
 \begin{equation*}
 \Gamma^{\ast}/\Gamma^+=\Gamma^{\ast}_p/\Gamma^+_p =
  (\Z/2\Z)^{\sharp\{\ell| d(B)\}}
 \end{equation*} 
\end{para}

\begin{para}\label{para:7.4.5} We now give some indications on Drinfeld's
 modular approach to \v{C}erednik's uniformization.  
\end{para}

\begin{lem}\label{lem:7.4.6}
 \begin{enumerate}
  \item Fake elliptic curves have potentially good reduction at every
	prime; the reduction is supersingular reduction at the critical
	prime $p$. In particular, over $\ovl{\F}_p$,
	they form a unique isogeny class.
  \item Fake elliptic curves have a unique principal $\ast$-polarization.
 \end{enumerate}
\end{lem}

\begin{proof}
 (i): due to the presence of many endomorphisms, the potential good
 reduction follows easily from the semi-stable reduction theorem for
 abelian schemes; also, the non-split quaternion algebra $B_p$ acts on
 the $p$-divisible group up to isogeny, and this forces the slopes to be
 $1/2$; for details and (ii), \cf \cite[III]{updcds}.
\end{proof}

 It follows from (i) that any fake elliptic curve over
 $\ovl{\F}_p$ is isogenous to the square of a supersingular elliptic
 curve $\ovl{E}$.  However, while there are only finitely many
 supersingular elliptic curves up to isomorphism, fake elliptic curves
 over $\ovl{\F}_p$ have continuous moduli! The moduli spaces are chains
 of $\P^1$'s.\footnote{Another curious feature of the higher
 dimensional case is that the product of $n > 1$ supersingular elliptic
 curves over $\ovl{\F}_p$ form a single isomorphism class (Deligne,
 Ogus)}

 Let us introduce the quaternion $\Q$-algebra $D$ ramified only
 at $p$ and $\infty$: $D\cong \End^\Q(\ovl{E})$. 
 Let $\ovl{A}$ be a fake elliptic curve over $\ovl{\F}_p$. It follows from
 \ref{lem:7.4.6} (i) that $\End^\Q(\ovl{A})\cong M_2(D)$, and $\End_{B}^\Q
 (\ovl{A})\cong \ovl{B}$.  

 In particular, the inner form $\ovl{G}$ of $G$ considered in \ref{sub:7.3} is
 $\ovl{B}^\times$, and $\ovl{G}(\Q_p)=GL_2(\Q_p)=J$.  

 It follows from \ref{lem:7.4.6} (ii) that we may disregard polarizations. The
 $p$-divisible groups attached to fake elliptic curves with
 $\mathcal{B}$-multiplication are special formal $\mathcal{B}_p$-modules
 of height $4$; recall from \ref{exa:4.3.4} the corresponding formal
 moduli space 
 $\displaystyle \mathfrak{M} = \coprod_\Z\;\mathfrak{M}^0$. The period
 mapping is an isomorphism (\ref{para:6.3.4}, \ref{sub:6.4}), and the
 commutative diagram of \ref{sub:7.3} specializes to

\begin{equation*}
 \begin{CD}
  \tilde{S}=\mathcal{M}^0 @>\mathcal{P}>> \mathcal{D}=\Omega_{\C_p} @.
  \subset\; \mathcal{D}^{\vee}=\P^1\\
  @V\mathcal{Q}VV  @VV\mathcal{Q}V @.\\
  S={\mathcal{X}^+(N)}_{\C_p}^{\mathrm{an}} @>\mathcal{P}>>
  \mathcal{D}/\Gamma^+_p(N)
 \end{CD}
\end{equation*}
 where the horizontal maps are isomorphisms (and the
 vertical maps are topological coverings if $\Gamma^+_p(N)$ is
 torsion-free), in complete analogy with \ref{sub:7.2}.
 
\begin{para}\label{para:7.4.7} Let us end this section by considering the
 case where {\it $p$ divides the level $N$}: $N=p^n.M,\; p\nmid M$. 

 In this case, $\Gamma^+_p(N)=\Gamma^+_p(M)$. The finite etale covering
 ${\mathcal{X}^+(N)}_{\C_p}^{\mathrm{an}}\rightarrow
 {\mathcal{X}^+(M)}_{\C_p}^{\mathrm{an}}$ is not a topological covering anymore
 (even if $\Gamma^+_p(M)$ is torsion-free). As was shown by Drinfeld, it
 is hidden in the finite etale covering $\Upsilon^n$ of $\Omega_{\C_p}$
 (\ref{para:6.3.5}).
 

 Let us assume for simplicity that $M>2$, so that $-1\notin
 \Gamma^+(M)$. This implies that $\Gamma^+(N)/\Gamma^+(M)\cong
 SL_1(\mathcal{B}/p^n\mathcal{B})$. Then $\Upsilon^n$ is a topological
 covering of ${\mathcal{X}^+(N)}_{\C_p}^{\mathrm{an}}$, and more precisely, one
 has the commutative diagram

 \begin{equation*}
 \begin{CD}
  {\mathcal{X}^+(N)}_{\C_p}^{\mathrm{an}}@>\sim>> \Upsilon^n/\Gamma^+_p(M)\\
  @VVV @VVV\\
  {\mathcal{X}^+(M)}_{\C_p}^{\mathrm{an}} @>\mathcal{P}>>
  \Omega_{\C_p}/\Gamma^+_p(M)
 \end{CD}
 \end{equation*}
 where the vertical maps are finite etale Galois
 covering maps with group $SL_1(\mathcal{B}/p^n\mathcal{B})$, \cf 
\cite[III.5.5]{updcds}, \cite{varshavsky98:_shimur}.  Similarly, there is such a $p$-adic global
 uniformization for any Shimura curve $\mathcal{X}_\Gamma$ (\cf
 \ref{para:7.4.1}).
\end{para}

\subsection{Other examples of global $p$-adic uniformization.}\label{sub:7.5}

\begin{para}\label{para:7.5.1}
 Shimura's paper \cite{shimura67:_const} deals not only with quaternion algebra
 over $\Q$, but also with quaternion algebras $B/E$ over a totally real
 number field $E$. When $B$ {\it splits at a single real place,} he
 constructs curves ${\mathcal{X}^+(N)}$ and ${\mathcal{X}^\ast}$
 satisfying properties similar to (a),(b),(c) of \ref{para:7.4.2}.


 In this context, $N$ may be an arbitrary ideal of $\mathcal{O}_E$,
 $\Q(\zeta_N)$ has to be replaced by the ray class field with conductor
 $N.\infty$, and the condition $N(g)>0$ means that the reduced norm of
 $g$ is totally positive, or, what amounts to the same, is positive at
 the split real place $\infty_0$.\index{000infinity0@$\infty_0$}
 
 The novelty, however, is that these Shimura curves have {\it no
 direct interpretation} as moduli spaces for
 decorated abelian varieties, since $B$ is not in the Albert list of
 endomorphism algebras of abelian varieties when $[E:\Q]>1$.


 To overcome this problem, Shimura introduces the following remarkable
 trick. Let $K$ be a totally imaginary quadratic extension of $E$. Then
 $B^\bullet =B\otimes_E K$ is in Albert's list (we endow it with the
 involution which is the tensor product of $\ast$ on $B$ and complex
 conjugation on $K$). It turns out that the Shimura varieties of PEL
 type for abelian varieties of dimension $g=4[E:\Q]$ ($\ast$-polarized,
 with level structure) with complex multiplication by $B^\bullet$ are
 twisted forms of the same curve. The Shimura curve attached to $B$ is
 obtained by Weil descent from the Shimura curves attached to
 $B\otimes_E K$, for sufficiently many quadratic extensions $K$.
\end{para}

\begin{para}\label{para:7.5.2}
 \v{C}erednik's original $p$-adic uniformization theorem
 already applies to Shimura curves attached to quaternion algebras over
 a totally real number field $E$. His method has been further developed
 in \cite{varshavsky98:_shimur}, \cite{varshavsky98:_shimur_ii}. On the
 other hand, Drinfeld's modular 
 approach has been extended by Boutot-Zink \cite{tpuosc}, using Shimura's
 trick. Let us sketch the principle of the construction.


 Let $G$ be the multiplicative group of $B^\times$, viewed as an
 algebraic group over $\Q$. Let $\Gamma$ be a congruence subgroup of
 $G^{\mathrm{ad}}(\Q)$ (alternatively, one can consider an open compact
 subgroup 
 $C$ of $G(\A_f)$). Let us first assume that $\Gamma$ is {\it maximal at
 $p$} (alternatively $C$ is of the form $C^pC_p$ for some open compact
 subgroup $C^p$ of $G(\A_f^p)$, \cf \ref{para:7.1.2}). The (non-connected) Shimura
 curve $\Sh$ given by
 \begin{equation*}
  \Sh_\C^{\mathrm{an}}= C\backslash \bigl((\P_\C^1\setminus\P^1(\R)) \times
   G(\A_f)\bigr) /G(\Q)
 \end{equation*}
 is defined over the reflex field, which is $E$ (embedded into $\R$ via
 $\infty_0$). Over some finite abelian extension $E'/E$, it decomposes as
 a disjoint union of geometrically connected Shimura curves
 $\mathcal{X}_\Gamma$ ($\mathfrak{h}/\Gamma$ over $\C\;$; one possible
 $\Gamma$ is the image of $G(\Q)\cap C$ in $\G^{\mathrm{ad}}(\Q)$).


 Let $v=v_0$ be a {\it critical} finite place of $E$, \ie such that
 $B_v$ is a non-split quaternion algebra. Let $p$ be the residue
 characteristic of $v_0$ and let $v_1,\ldots, v_m$ be the other places
 of $E$ above $p$.


 Let $K$ be a totally imaginary quadratic extension of $E$, such that
 every place $v_i$ splits in $K$. We denote by $w_i, \ovl{w_i}$ the two
 places of $K$ above $v_i$. The quaternion algebra $B^\bullet =
 B\otimes_E K$ is ramified at $w=w_0$: $B^\bullet_w\cong B_v$.


 We fix an extension $\infty_0: K\hookrightarrow \C$ of the real
 embedding $\infty_0$ of $E$, and a double embedding $\C \hookleftarrow
 \ovl{\Q} \hookrightarrow \C_p$ whose restriction to $K$ is compatible
 with $(\infty_0,w)$. The embeddings $K \hookrightarrow \C_p$ which
 factor through some $w_i$ but not through some $\ovl{w_i}$ form a CM
 type $\Psi$: $\Psi\coprod \ovl{\Psi}=\Hom(K,\C_p)$ (identified with
 $\Hom(K,\C)$ through our chosen double embedding).
 

There is a natural $\Q$-group $G^\bullet$ such that
 $G^\bullet(\Q)=\{b\in B^\bullet| bb^\ast\in E\}$, acting on
 $V^\bullet=V\otimes_E K$. One can attach to $C$ an open compact
 subgroup $C^\bullet$ of $G^\bullet(\A_f)$, maximal at $p$, and a
 Shimura curve $\Sh^\bullet$ such that $\Sh$ is open and closed in
 $\Sh^\bullet$ (\cite[3.11]{tpuosc}). Alternatively, one can attach to
 $\Gamma$ a congruence subgroup $\Gamma^\bullet$ in $G^{\bullet ad}(\Q)$
 and a corresponding connected Shimura curve
 $\mathcal{X}_{\Gamma^\bullet}$ which is a twisted form of
 $\mathcal{X}_{\Gamma}$ (defined over some abelian extension of
 $K$). The point is that this new Shimura curve $\Sh^\bullet$ is of PEL
 type. More precisely, for $C$ small enough, it is a moduli space for
 $\ast$-polarized abelian varieties of dimension $g=4[E:\Q]$ with
 multiplication by
 $\mathcal{B}^\bullet=\mathcal{B}\otimes_{\mathcal{O}_E}\mathcal{O}_K$,
 Shimura type $\Psi$, and level structure. As in \ref{pro:7.1.5}, it
 admits a model $\Sh^\bullet$ over $\mathcal{O}_{E_v}$.


 Let $\udl{A} \rightarrow \Sh^\bullet$ be the universal abelian
scheme up to prime to $p$ isogeny.  Due to the $\mathcal{O}_K$-action,
the $p$-divisible group of $\udl{A} $ splits as $
\udl{A}[p^\infty]\cong \prod_i \Lambda_{w_i}\times \Lambda_{\ovl{w_i}}$,
where $ \Lambda_{w_i}$ and $\Lambda_{\ovl{w_i}}$ are $p$-divisible groups
of height $4[E_{v_i}:\Q_p]$ with formal multiplication by
$E_{v_i}$. Moreover, $ \Lambda_{w}= \Lambda_{w_0}$ is a special formal
$\mathcal{B}_v$-module of height $4.[E_{v}:\Q_p]$ (defined as in 
\ref{exa:4.3.4}, 
replacing $\Q_p$ by $E_v$), and $ \Lambda_{\ovl{w}}$ is the Serre dual of
$\Lambda_{w}$; the other $ \Lambda_{w_i},\;i>0,$ are etale, and the $
\Lambda_{\ovl{w_i}}$ are their respective Serre duals. Hence the Newton
polygon of the Dieudonn\'e module of $\udl{A}[p^\infty]$ has slopes
$0,\frac{1}{2},1$ (only $\frac{1}{2}$ if $m=0$).

\vspace*{2mm}
%#BEGIN
\begin{figure}[h]
 \begin{picture}(120,70)(-20,-10)
  \put(-10,0){\line(1,0){110}}
  \put(0,-10){\line(0,1){70}}
  \bezier{10}(60,15)(60,7.5)(60,0)
  \bezier{30}(90,45)(90,22.5)(90,0)
  \put(15,5){\makebox(0,0)[b]{{\tiny $0$}}}
  \put(40,15){\makebox(0,0)[c]{{\tiny $1/2$}}}
  \put(68,32){\makebox(0,0)[c]{{\tiny $1$}}}
  \thicklines
  \put(0,0){\line(1,0){30}}
  \put(30,0){\line(2,1){30}}
  \put(60,15){\line(1,1){30}}
 \end{picture}
 \caption{}
\end{figure}
%#END

Let $\mathfrak{v}$ be the discrete valuation ring
$\widehat{\Z^{\mathrm{ur}}}.\mathcal{O}_{E_v}$ as in
\ref{para:7.1.2}. The moduli problem 
for $\mathcal{B}_v$-module of height $4.[E_{v}:\Q_p]$ is representable
by the $p$-adic formal scheme $\displaystyle \coprod_\Z
\;({}_{E_v}\hat\Omega)_{\mathfrak{v}}$, where ${}_{E_v}\hat\Omega$ is
the analogue for (and over) $\mathcal{O}_{E_v}$ of the formal scheme
$\hat\Omega$, with associated analytic space
$(_{E_v}\Omega)_{\C_p}=\P^1_{\C_p}\setminus \P^1(E_v)$, \cf \ref{para:6.3.5}.
\end{para}

\begin{para}\label{para:7.5.3}
 This is the modular way taken by Boutot and Zink to reprove (and
 strengthen) \v{C}erednik's theorem.  As was remarked by
 Varshavsky \cite[3.13]{varshavsky98:_shimur_ii}, one obtains a simpler adelic
 formulation if one works with a twisted version ${}^{\mathrm{ad}}\Sh$ of $\Sh$
 (corresponding to a different normalization of $h_0$ in \ref{para:1.2.4}); this
 does not change the adelic description of the associated complex space
 \begin{equation*}
  {}^{\mathrm{ad}}\Sh_\C^{\mathrm{an}}= C\backslash
   \bigl((\P_\C^1\setminus\P^1(\R)) \times G(\A_f)\bigr)/G(\Q),
 \end{equation*}
 but makes the analogy with the following description of the associated
 $p$-adic space more striking.
 

 Let $\ovl{B}$ be the quaternion algebra over $E$ obtained from $B$ by
 interchanging the invariants at $\infty_0, v$, and let $\ovl{G}$ be the
 $\Q$-algebraic group attached to $\ovl{B}^\times$. Let us identify
 $\ovl{G}(\A^p_f)$ with $G(\A^p_f)$ via an anti-isomorphism
 $\ovl{B}\otimes_E \A^p_f \rightarrow B\otimes_E \A^p_f ,$ and set
 $\ovl{C}=GL_2(E_v)C^p$ viewed as a subgroup of $\ovl{G}(\A_f)$.  Then
 \begin{align*}
  {}^{\mathrm{ad}}\Sh_{\C_p}^{\mathrm{an}} & =C^p \backslash
  \bigl((_{E_v}\Omega)_{\C_p}\times G(\A_f^p)\bigr)/{\ovl{G}}(\Q)\\
  & =\ovl{C}\backslash \bigl((\P^1_{\C_p}\setminus\P^1(E_v))\times
  \ovl{G}(\A_f)\bigr)/{\ovl{G}}(\Q)
 \end{align*}
 where ${\ovl{G}}(\Q)$ acts on $(_{E_v}\Omega)_{\C_p}$ through
 ${\ovl{G}}^{\mathrm{ad}}(\Q_p)\cong PGL_2(E_v)$.  

 Moreover, one can drop the assumption that $C$ is maximal at $p$, on
 replacing $(_{E_v}\Omega)_{\C_p}$ by some finite etale covering as in
 \ref{para:7.4.7}. \cite[3]{tpuosc}, \cite[5]{varshavsky98:_shimur_ii}.
\end{para}

\begin{para}\label{para:7.5.4}
 Rapoport, Zink and Boutot on one hand (by modular methods \`a la
 Drinfeld), and Varshavsky on the other hand (by group-theoretic methods
 \`a la \v{C}erednik), have pointed out similar global
 $p$-adic uniformizations in much more general cases.

 Let us only mention a remarkable example: Mumford's {\it fake
 projective plane} \cite{aaswkak9pq0}. This is a smooth projective complex
 surface $S$ with the same Betti numbers as $\mathbb{P}^2$, but not
 isomorphic to $\mathbb{P}^2$. Actually, it is a surface of general
 type. Mumford's construction is of diadic nature. It is a quotient of
 the two-dimensional generalization $\Omega^{(2)}$ of $\Omega$ by an
 explicit arithmetic group. This $\Omega^{(2)}$ is the complement of the
 lines defined over $\Q_2$ in the projective plane.

 Using the last theorem of \cite{psfpg} or \cite{varshavsky98:_shimur_ii}, one can see
 that $S$ is actually a Shimura surface of PEL type, parametrizing
 polarized abelian varieties of dimension $9$, with multiplication by an
 order in a simple $\Q$-algebra $B$ of dimension $18$ with center
 $\Q(\sqrt{-7})$. The invariants of $B$ at the two primes of
 $\Q(\sqrt{-7})$ above $2$ are $1/3$ and $2/3$ (\cf a recent
 preprint by F. Kato for more detail). The associated group $G$ is a
 group of unitary similitudes of signature $(1,2)$.


 The ``diadic ball'' $\Omega^{(2)}$ (or rather its formal avatar) has an
 modular interpretation for certain $p$-divisible groups (with
 multiplication by $B_2$). These $p$-divisible groups split as a product
 of three $p$-divisible groups of slope $1/3$ and height $3$, and three
 $p$-divisible groups of slope $2/3$ and height $3$.


 The associated Gauss-Manin connection splits into six factors of rank
 $3$ (in two variables). It resembles the Appell-Lauricella connection
 studied in \cite{pdrefapdfhdpv}, \cite{calip}, \cite{fde}, but is different.
\end{para}

\subsection{Application of the theory of $p$-adic Betti lattices.}
\label{sub:7.6}

\begin{para}\label{para:7.6.1}
 We consider more closely the Gauss-Manin connection
 $(\mathcal{H},\nabla_{\mathrm{GM}})$ in the case of Shimura curves.


 We begin with the case $E=\Q$ as in \ref{sub:7.4}.  Then $\mathcal{H}$
 is a vector bundle of rank $4$, and the $B$-action on $\mathcal{H}$
 obtained by functoriality of De Rham cohomology commutes with
 $\nabla_{\mathrm{GM}}$.

 Let us choose an imaginary quadratic field $F=\Q(\sqrt{-d})$ inside
 $B$. Then $(\mathcal{H},\nabla_{\mathrm{GM}})\otimes F$ splits into two parts:
 $(\mathcal{H},\nabla_{\mathrm{GM}})_+\oplus
 (\mathcal{H},\nabla_{\mathrm{GM}})_{-}$ (on 
 which $F\subset B$ acts through identity and complex conjugation
 respectively).


 On the other hand, let us fix embeddings $F\subset \ovl{\Q} \subset \C$,
 and let us fix a base point $s\in \mathcal{X}_{\Gamma}(\ovl{\Q})$ of our
 Shimura curve.  This is the moduli point of a decorated fake elliptic
 curve $A/\ovl{\Q}$. We have $V^\C=\mathrm{H}^1_B(A_\C,\C)\cong
 (\mathcal{H}_\C^{\mathrm{an}})^{\nabla_{\mathrm{GM}}}$, and this again splits
 into $\pm$ parts. The Betti cohomology $F$-space
 $\mathrm{H}^1_B(A_\C,F)$ already splits into two parts:
 $\mathrm{H}^1_B(A_\C,F)_+\oplus \mathrm{H}^1_B(A_\C,F)_{-}$, both of
 rank two and stable under the action of $B\otimes F \cong M_2(F)$.

 The flag space $\mathcal{D}^\vee$ is actually defined over the reflex
 field $\Q$. Over $F$, it becomes isomorphic to $\P^1$, or more
 precisely, to $\P(\mathrm{H}^1_B(A_\C,F)_+)$. The period mapping describes the
 ``slope'' of the Hodge line $F^1\cap V^\C_+$ with respect to a basis of
 $\mathrm{H}^1_B(A_\C,F)_+\subset
 ((\mathcal{H}_+)_\C^{\mathrm{an}})^{\nabla_{\mathrm{GM}}}$, and is 
 thus given by a quotient of solutions of the (partial) Gauss-Manin
 connection $\nabla_{\mathrm{GM}+}$.


 Assume for simplicity that $\Gamma\subset B^\times/\Q^\times$ is
 torsion-free. Then $\Gamma$ may be identified with the projective
 monodromy group, \ie with the quotient by the homotheties of the image
 of the monodromy representation
 \begin{equation*}
  \pi_1((\mathcal{X}_{\Gamma })_\C,s)\rightarrow GL(\mathrm{H}^1_B(A_\C,F)_+).
 \end{equation*}


\end{para}

\begin{para}\label{para:7.6.2} Let us now turn to the $p$-adic side. We fix an
embedding $\ovl{\Q} 
\subset \C_p$. We shall limit ourselves to the case of a prime $p$ of
 {\it supersingular reduction} for $A$. Note that this includes the case
 of critical primes (\ref{lem:7.4.6}).


 Let $D$ be the definite quaternion algebra over $\Q$ which is ramified
 only at $p$, and let $\ovl{B}$ be the definite quaternion algebra over
 $\Q$ which is ramified at the same primes as $B$ except $p$. Then,
 whether $p$ is critical or not, the group $J$ (of self-quasi-isogenies
 of $\ovl{A}[p^\infty]$ commuting with $B$) is $\ovl{B}_p^\times$.
 

 By the local criterion for splitting fields \cite[III.3.5]{adadq}, $F$
 is also a splitting field for the quaternion algebra $D$ (and $\ovl{B}$
 as well). We can therefore apply the theory of
 I.\ref{sub-p-adic-betti-lattice-supersing}, and get an 
 $F$-structure $\mathrm{H}^1_B(A_{\C_p},F)$ inside
 $\mathrm{H}^1_{\mathrm{cris}}(\ovl{A},/\C_p)\cong
 V^{\C_p}=\D(\ovl{\Lambda})\otimes \C_p.$ By construction, it is stable
 under the action of $M_2(D)\otimes F\cong M_4(F)\cong B\otimes \ovl{B}
 \otimes F$.


 Using the embedding $F\otimes 1\otimes F\subset B\otimes \ovl{B} \otimes
 F$, one splits this $F$-structure $\mathrm{H}^1_B(A_{\C_p},F)$ into two parts
 $\mathrm{H}^1_B(A_{\C_p},F)_+\oplus \mathrm{H}^1_B(A_{\C_p},F)_{-}$ (on
 which $F\otimes 
 1\otimes 1\subset B\otimes \ovl{B} \otimes F$ acts through identity and
 complex conjugation respectively), both of rank two and stable under
 the action of $\ovl{B}\otimes F \subset M_4(F)$, hence under $J$.  \par
 The $p$-adic flag space $\mathcal{D}^\vee$ is again defined over the
 reflex field $\Q$. Over $F$, it becomes isomorphic to
 $\P(\mathrm{H}^1_B(A_{\C_p},F)_+)$ (note that the scaling ambiguity of
 $\mathrm{H}^1_B(A_{\C_p},F)_+$ by a factor in $\sqrt{F^\times}$ disappears
 here). The $J$-equivariant period mapping $\mathcal{P}$ describes the
 ``slope'' of the Hodge line $F^1\cap V^{\C_p}_+$ with respect to a
 fixed basis $(e_0,\,e_1)$ of $\mathrm{H}^1_B(A_{\C_p},F)_+\subset
 ((\mathcal{H}_+)_{\C_p}^{\mathrm{an}})^{\nabla_{\mathrm{GM}}}$, and is
 thus given by a quotient of solutions of $\nabla_{\mathrm{GM}+}$.


 Note that if $B$ is split at $p$, then $\mathcal{P}$ is of
 Gross-Hopkins type, and one has only a local $p$-adic uniformization of
 $\mathcal{X}_{\Gamma}$ around $s$ (in a supersingular disk).


 If instead $p$ is critical, then $\mathcal{P}$ is of Drinfeld type, and
 the group $\Gamma_p \subset J^1$ appearing in the global $p$-adic
 uniformization of $\mathcal{X}_{\Gamma}$ may be identified with the
 quotient by the homotheties of the image of the monodromy representation
 (in the sense of \ref{cor:7.2.6})
 \begin{equation*}
  \pi_1^{\mathrm{top}}((\mathcal{X}_{\Gamma })_{\C_p},s)\rightarrow
   GL(\mathrm{H}^1_B(A_{\C_p},F)_+).
 \end{equation*}

 A slight refinement provides a rank two
 $\mathcal{O}_F[\frac{1}{p}]$-lattice
 $\mathrm{H}^1_B(A_{\C_p},\mathcal{O}_F[\frac{1}{p}])_+$ in
 $((\mathcal{H}_+)_{\C_p}^{\mathrm{an}})^{\nabla_{\mathrm{GM}}}$ stable
 under $p$-adic monodromy.
\end{para}

\begin{para}\label{para:7.6.3}
 One can go one step further and see that the phenomenon encountered in
 \ref{para:6.3.3} shows up here again: ``the'' quotient $\tau$ of solutions of
 solutions of ${\nabla_{\mathrm{GM}}}$ which expresses locally the
 period mapping 
 $\mathcal{P}$ is essentially given by the same formula in the complex
 and the $p$-adic cases.


More precisely, assume that $A$ has complex multiplication by an order
in $M_2(F)$ (so that $A$ is isogenous to the square of an elliptic curve
with complex multiplication by $\mathcal{O}_{\Q(\sqrt{-d})}$ with
[fundamental] discriminant $-d$ ), and identify $\End_B(A)$ with $F$. We
denote by $\epsilon$ the Dirichlet character, $w$ the number of roots of
unity, $h$ the class number attached to $\mathcal{O}_{\Q(\sqrt{-d})}$.


\noindent We may choose our symplectic basis $e_0,e_1$ of
$\mathrm{H}^1_B(A_{\C},F)_+$ (\resp $\mathrm{H}^1_B(A_{\C_p},F)_+$) in
 such a way that 
$[\sqrt{-d}]^\ast e_0=-\sqrt{-d} \; e_0,\,$ $[\sqrt{-d}]^\ast
e_1=\sqrt{-d} \;e_1$ (here $[\sqrt{-d}]$ is viewed as an element of
$\End_B(A)=F$).

\noindent Let $\omega$ be a section of $\mathcal{H}_+$ in a Zariski
neighborhood $U$ of $s$, which corresponds to a relative differential of
the first kind (no pole) on the universal abelian scheme. Then $\omega$
is a cyclic vector for $\mathcal{H}_+$ over $U$ with respect to
$\nabla_{\mathrm{GM}}$, so that $\nabla_{\mathrm{GM}}$ amounts to a
 concrete linear 
differential equation of order two over $U$.

\noindent We calculate $\tau$ with respect the basis $e_0, e_1$
 of $\mathrm{H}^1_B(A_{\C},F)_+$ (\resp $\mathrm{H}^1_B(A_{\C_p},F)_+$)
 in a small 
 neighborhood of the base point $s$. Denoting by $z$ an algebraic local
 parameter at $s$ in $U$, and extending $e_0,\, e_1$ by horizontality, we
 then have the proportionality
\begin{equation*}
 \omega \sim \tau(z) e_0 + e_1
\end{equation*}
\end{para}

\begin{thm}\label{thm:7.6.4}
 There is a quotient $y=y_1/y_2$ of solutions of $\nabla_{\mathrm{GM}}$ in
 $\ovl{\Q}[[z]]$ such that $\tau$ takes the form
 $\;\tau(z)=\kappa^{-1}\;y(z)\;$ for a suitable constant $\kappa$.

\noindent In the complex case, $\;\kappa\sim\displaystyle
 \prod_{u\in (\Z/d)^\times}
 \Bigl(\Gamma\Bigl\langle\frac{u}{d}\Bigr\rangle\Bigr)^{\epsilon(u)w/2h}\;\in
 \C^\times/\ovl{\Q}^\times$.

\noindent In the $p$-adic case, assuming moreover that $p\nmid d $, 
\begin{equation*}
 \kappa\sim \prod_{u\in (\Z/d)^\times}
  \Bigl(\Gamma_p\Bigl\langle\frac{pu}{d}\Bigr\rangle\Bigr)^{-\epsilon(u)w/4h}
  \in \C_p^\times/\ovl{\Q}^{\times}.
\end{equation*}
\end{thm}
\begin{proof}
 Let us write $\omega_{| z=0}=\omega_{11}e_1$. Then
 $\tau(z)=\kappa^{-1} y(z)$, with $\kappa = \omega_{11}^2/2\pi i$ in
 the complex case, $\kappa = \omega_{11}^2$ in the $p$-adic case (the
 factor $2\pi i$ which arises in the complex case is the factor of
 proportionality of the symplectic forms in De Rham and in Betti
 cohomology respectively). The evaluation of $\omega_{11}$ in the
 complex case is given by the Lerch-Chowla-Selberg formula (\cf
 I.\ref{thm-ogus-formula}). In the $p$-adic case, it is given by
 I.\ref{thm-cm-period}.
\end{proof}

\begin{para}\label{para:7.6.5}
 This discussion generalizes without much complication to the case of
 Shimura curves attached to a quaternion algebra $D$ over a totally real
 number field $E$ of degree $>1$. We concentrate on the case of a
 critical place $v$ and assume, for simplicity, that $E$ is Galois over
 $\Q$ and that $v$ is the unique place of $E$ with the same residue
 characteristic $p$ (\ie $m=0$ in the notation of \ref{para:7.5.2}).


 Let $F$ be a totally imaginary quadratic extension of $E$ contained in
 $B$. Then $F$ is a splitting field for $B,\ovl{B}$, and $D$. Let us fix
 a double embedding $\C \hookleftarrow \ovl{\Q} \hookrightarrow \C_p$ as
 in \ref{para:7.5.2}, and an extension $FK\hookrightarrow \ovl{\Q}$ of $\infty_0$.
 For $C$ small enough, $\Sh^\bullet$ carries a universal abelian scheme
 of relative dimension $g=[FK:\Q]$. After tensoring $\otimes_K
 \ovl{\Q}$, its Gauss-Manin connection splits into pieces of rank two,
 indexed by the embeddings of $FK$ into $\ovl{\Q}$. We select the piece
 $(\mathcal{H},\nabla_{\mathrm{GM}})_+$ corresponding to $\infty_0$.


 Let us fix a base point $s\in \Sh(\ovl{\Q})\subset \Sh^\bullet(\ovl{\Q})
 $ of our Shimura curve.  This is the moduli point of a decorated
 abelian variety $A^\bullet/\ovl{\Q}$ (of dimension $g=[FK:\Q]$ and
 multiplication by $\mathcal{B}^\bullet$). We have $V^\bullet=V\otimes_E
 K \cong \mathrm{H}^1_B(A^\bullet_\C,\Q)$.  Note that
 $\mathrm{H}^1_B(A_\C,F)\cong 
 V\otimes_\Q FK$ splits into pieces of $FK$-rank two, indexed by the
 embeddings of $F$ into $\ovl{\Q}$. We select the piece
 $\mathrm{H}^1_B(A^\bullet_\C,F)_+$ corresponding to $\infty_0$.  We
 then have a 
 natural embedding $\mathrm{H}^1_B(A^\bullet_\C,F)_+\subset
 ((\mathcal{H}_+)_{\C}^{\mathrm{an}})^{\nabla_{\mathrm{GM}}}$, and
 $\mathrm{H}^1_B(A^\bullet_\C,F)_+$ is stable under monodromy.


 On the other hand, it follows from \ref{lem:7.4.6}, that $A^\bullet$ has
 supersingular reduction $\ovl{A}^\bullet$ at $v$. We can construct the
 $F$-structure $\mathrm{H}^1_B(A^\bullet_{\C_p},F)$ inside
 $\mathrm{H}^1_{\mathrm{cris}}(\ovl{A}^\bullet,/\C_p)\cong
 V^{\C_p}=\D(\ovl{\Lambda}^\bullet)\otimes \C_p.$ 
 By construction, it is stable under the action of $M_g(D)\otimes F\cong
 M_{2g}(F)$, which contains naturally $B\otimes_\Q \ovl{B}\otimes_\Q F$.


 Using the embedding $F\otimes 1\otimes F\subset B\otimes \ovl{B} \otimes
 F$, one splits this $F$-structure $\mathrm{H}^1_B(A_{\C_p},F)$ into parts of
 $FK$-rank two, indexed by the embeddings of $F$ into $\ovl{\Q}$. We
 select the piece $\mathrm{H}^1_B(A^\bullet_{\C_p},F)_+$ corresponding to
 $\infty_0$.  We then have a natural embedding
 $\mathrm{H}^1_B(A^\bullet_{\C_p},F)_+\subset
 ((\mathcal{H}_+)_{\C_p}^{\mathrm{an}})^{\nabla_{\mathrm{GM}}}$, and
 $\mathrm{H}^1_B(A^\bullet_{\C_p},F)_+$ is stable under $p$-adic monodromy.
\end{para}

\subsection{Conclusion.}\label{sub:7.7}

 Roughly speaking, the period mapping $\mathcal{P}$ attaches to
each member of an algebraic family of algebraic complex varieties (with
additional structure) a point in some flag space, which encodes the
Hodge filtration in the cohomology of this variety. This mapping can be
expressed in terms of quotients of solutions of the fuchsian
differential equation (Picard-Fuchs/Gauss-Manin) which controls the
variation of the cohomology.


We have described the theory of period mappings for $p$-divisible groups
in somewhat similar terms.

\bigskip
 Sometimes, it is possible to go through the looking-glass of sheer
 analogies. This occurs when we 
 restrict our attention to $p$-divisible groups attached to polarized
 abelian varieties with prescribed endomorphisms and fixed ``Shimura
 type''. In this situation, we have three algebraic objects defined over
 a number field (the reflex field):

\begin{enumerate}
 \renewcommand{\theenumi}{\alph{enumi}}
 \item the Shimura variety $\Sh$ parametrizing certain
       ``decorated'' abelian varieties,
 \item the Gauss-Manin connection $\nabla_{\mathrm{GM}}$ describing
       the variation of these abelian varieties over $\Sh$, and
 \item the flag variety $\mathcal{D}^\vee$.
\end{enumerate}


 These objects are transcendentally related, both in the complex and
 $p$-adic sense.

 The period mapping may be viewed as
 multivalued and defined on $p$-adic open domains in $\Sh$; it is given
 as in the complex case by quotients of solutions of the Gauss-Manin
 connection (as a $p$-adic connection).


 In most situations, these open domains are small and $\mathcal{P}$ is
 single-valued.


 However, it happens in some remarkable cases that $\mathcal{P}$ is a
 global multivalued function. This occurs when the abelian varieties
 under consideration form a {\it single isogeny class modulo $p$}. In
 this situation, the complex and $p$-adic transcendental relations
 between $\Sh, \nabla_{\mathrm{GM}}$ are completely similar: $\Sh$ is
 uniformized 
 by an analytic space which is etale over some open subset $\mathcal{D}$
 of $\mathcal{D}^{\vee}$, and the group of deck transformations $\Gamma$
 is isomorphic to the projective global monodromy group of
 $\nabla_{\mathrm{GM}}$.  In III, we shall analyze in detail the
 group-theoretic 
 aspects of this situation.
